% !TEX program = xelatex
% ============================================================
% 玉器门店顾客服务与经营管理系统 —— 业务建模、需求、分析与设计
% 课程：系统分析与设计    编译：xelatex（ctexart，中文必须用 xelatex）
% C 阶段全部分析图由 docs/puml/分析图/ 渲染到 docs/figs/analysis/：
% 05a/05a2 边界类与控制类、05b1-05b16 实体类关联图逐张；
% D 阶段的设计图由 docs/puml/设计图/ 渲染到 docs/figs/design/。
% ============================================================
\documentclass[12pt, a4paper]{ctexart}

\usepackage[top=2.4cm, bottom=2.4cm, left=2.6cm, right=2.6cm]{geometry}
\usepackage{graphicx}
\usepackage{booktabs}
\usepackage{longtable}
\usepackage{array}
\usepackage{tabularx}
\usepackage{enumitem}
\usepackage{xcolor}
\usepackage{caption}
\usepackage{float}
\usepackage{needspace}
\usepackage{pdflscape}
\usepackage{hyperref}

\hypersetup{colorlinks=true, linkcolor=black, urlcolor=blue, citecolor=black}
\captionsetup{font=small, labelfont=bf}

% 图片占位命令：依次在 figs/ 的 analysis/design 两个分类子文件夹里找文件，
% 都没有时给出提示方框，保证文档随时可编译；文件名做 label
\newlength{\figw}
\newcommand{\fig}[3][1.0]{%
  \setlength{\figw}{#1\textwidth}%
  \clearpage
  \begin{figure}[H]
    \centering
    \IfFileExists{figs/analysis/#3.png}{%
      \includegraphics[width=\figw,height=20cm,keepaspectratio]{figs/analysis/#3.png}%
    }{%
    \IfFileExists{figs/analysis/#3}{%
      \includegraphics[width=\figw,height=20cm,keepaspectratio]{figs/analysis/#3}%
    }{%
    \IfFileExists{figs/design/#3.png}{%
      \includegraphics[width=\figw,height=20cm,keepaspectratio]{figs/design/#3.png}%
    }{%
    \IfFileExists{figs/design/#3}{%
      \includegraphics[width=\figw,height=20cm,keepaspectratio]{figs/design/#3}%
    }{%
      \fbox{\parbox[c][4.6cm][c]{0.8\textwidth}{\centering
        【待插入图形】\vspace{6pt}
        \texttt{figs/#3}}}}}}}
    \caption{#2}
    \label{fig:#3}
  \end{figure}%
}

\newcommand{\figtwo}[5][1.0]{%
  \setlength{\figw}{#1\textwidth}%
  \clearpage
  \begin{figure}[H]
    \centering
    \IfFileExists{figs/analysis/#3.png}{\includegraphics[width=\figw]{figs/analysis/#3.png}}{%
    \IfFileExists{figs/design/#3.png}{\includegraphics[width=\figw]{figs/design/#3.png}}{\fbox{【待插入图形】\texttt{#3}}}}%
    \caption{#2}\label{fig:#3}\vspace{6pt}
    \IfFileExists{figs/analysis/#5.png}{\includegraphics[width=\figw]{figs/analysis/#5.png}}{%
    \IfFileExists{figs/design/#5.png}{\includegraphics[width=\figw]{figs/design/#5.png}}{\fbox{【待插入图形】\texttt{#5}}}}%
    \caption{#4}\label{fig:#5}
  \end{figure}%
}

% 同页两图、各自独立宽度：#1=宽1 #2=图题1 #3=图名1 #4=宽2 #5=图题2 #6=图名2
\newcommand{\figtwoW}[6]{%
  \clearpage
  \begin{figure}[H]
    \centering
    \IfFileExists{figs/analysis/#3.png}{\includegraphics[width=#1\textwidth]{figs/analysis/#3.png}}{%
    \IfFileExists{figs/design/#3.png}{\includegraphics[width=#1\textwidth]{figs/design/#3.png}}{\fbox{【待插入图形】\texttt{#3}}}}%
    \caption{#2}\label{fig:#3}\vspace{6pt}
    \IfFileExists{figs/analysis/#6.png}{\includegraphics[width=#4\textwidth]{figs/analysis/#6.png}}{%
    \IfFileExists{figs/design/#6.png}{\includegraphics[width=#4\textwidth]{figs/design/#6.png}}{\fbox{【待插入图形】\texttt{#6}}}}%
    \caption{#5}\label{fig:#6}
  \end{figure}%
}
\usepackage{caption}
\captionsetup{labelsep=quad}
\counterwithin{table}{section}
\counterwithin{figure}{section}
\renewcommand{\thetable}{\thesection-\arabic{table}}
\renewcommand{\thefigure}{\thesection-\arabic{figure}}
\newcommand{\code}[1]{\texttt{\small #1}}
% 带圈数字划归 CJK 类，由中文字体渲染，避免静默丢字
\xeCJKDeclareCharClass{CJK}{"2460 -> "24FF, "2014}

\title{玉器门店顾客服务与经营管理系统\\[6pt] \large 愿景、业务建模、需求、分析与设计}
\author{姓名：\underline{张昊轩}\quad 学号：\underline{8102240429}\quad 班级：\underline{信管 T2401}}
\date{课程：系统分析与设计 \quad 指导教师：毕文杰 \quad 2026 年 10 月}

\newcommand{\figab}[3][0.86]{%
  \needspace{10\baselineskip}
  \begin{figure}[H]
    \centering
    \IfFileExists{figs/business-use-cases/#3}{%
      \includegraphics[width=#1\textwidth]{figs/business-use-cases/#3}%
    }{%
    \IfFileExists{figs/business-sequence/#3}{%
      \includegraphics[width=#1\textwidth]{figs/business-sequence/#3}%
    }{%
    \IfFileExists{figs/improved-business-sequence/#3}{%
      \includegraphics[width=#1\textwidth]{figs/improved-business-sequence/#3}%
    }{%
    \IfFileExists{figs/system-use-cases/#3}{%
      \includegraphics[width=#1\textwidth]{figs/system-use-cases/#3}%
    }{%
      \fbox{\parbox[c][4.6cm][c]{0.8\textwidth}{\centering
        【待插入图形】\\[6pt]
        \texttt{figs/#3}}}}}}}
    \caption{#2}
    \label{fig:#3}
  \end{figure}%
}

\begin{document}
\maketitle

\tableofcontents

\newpage

\section{A 业务建模}
% ============================================================

\subsection{案例背景}

老板（柜台）在商场经营一个玉器门店（档口），在柜 4–5 人，销售手镯、吊坠、摆件等成品玉器。
玉器经营与普通零售不同：每件玉器都是独一件、一物一价，客单价从一两千元到六位数不等；
在柜货品约 3000–4000 件（小件多、部分货品可拆，数量常有出入），货值约 2000–3000 万元，
其中相当一部分是向同行调来的货，并非都是门店自有。
货品来源按占比大致是：同行调货约 40\%、买料加工约 30\%、回收回流约 20\%、成品进货约 10\%；
主做新疆料，加工以切山料、籽料为主，戈壁料与籽料也卖一些，原石山流水料因风险大而少买。
正因为非标准化——同一款手串做 100 条，克重、成色、定价各不相同，无法按“件”套用现成进销存，
只能逐件建档。玉的真假依靠鉴定证书，而产地目前证书做不出来，靠老板眼力把关。

随着经营规模扩大，管理方式没有相应跟上，问题集中在四处：

\begin{enumerate}[leftmargin=1.6em, itemsep=2pt]
  \item 调货往来靠微信：调货时在图片上写价、发图给同行，货拿走就算成交，
        微信记录会丢、图片会被清理，事后认不出调出去的是哪一件，对账经常扯皮——这是最头疼的一处；
  \item 进销记录靠纸质台账和记忆：每件货品从入库到售出没有系统记录，
        成本和入库时间要翻进货单据回忆，货品周转周期算不出来，压柜的货越来越多；
  \item 客户信息散落：客户联系方式、购买记录、定金与改圈口等约定散落在微信聊天和手写小票里，
        想做精准营销没有数据可用，月度销售客户数要月底翻账本数出来；
  \item 经营分析缺数据：销售额靠月底翻账本数出来，按货品、客群的分析无法进行。
\end{enumerate}

为此决定开发玉器门店顾客服务与经营管理系统，把每件货品从进货入库到销售回款的全程记录起来。

\subsection{涉众与期望}

\needspace{8\baselineskip}
\begin{table}[H]
\centering
\caption{涉众及其利益}
\begin{tabularx}{\textwidth}{@{}l l X@{}}
\toprule
涉众 & 与业务的关系 & 涉众利益 \\
\midrule
门店老板 & 业务工人、经营决策者 & 掌握每件货品的进销存与周转周期，减少积压；月度销售额与客户数随手可查；回收、调货往来不错账，出错了追得到经手人 \\
销售 & 业务工人、一线执行者 & 记录过程简单、不干扰成交；成交归属清楚；老板不在店时代理收货与送检也有据可查 \\
客户 & 业务执行者 & 交易顺畅，价格清楚，证书可信，售后（退货）有据，鉴赏期内能从容决定 \\
玉料供应商、成品玉器供应商 & 业务执行者 & 供货与结算账目清楚 \\
加工师傅 & 业务执行者 & 加工要求、交期与加工费明确，加工过程可回看 \\
鉴定机构 & 业务执行者 & 送检样品与证书信息完整 \\
回流出让方 & 业务执行者 & 回收作价合理、付款及时 \\
同行商户 & 业务执行者 & 调货往来账目清楚，不再靠微信图片对账 \\
\bottomrule
\end{tabularx}
\end{table}

\noindent
用例规约中的每条约束都回溯到上表中某条利益，见 B 部分各用例规约的涉众利益与补充约束栏。

\subsection{愿景}

愿景写的是引进系统后目标组织要改进的指标。

\needspace{8\baselineskip}
\begin{table}[H]
\centering
\caption{愿景（目标系统：玉器门店顾客服务与经营管理系统）}
\begin{tabularx}{\textwidth}{@{}p{2.6cm} X@{}}
\toprule
项目 & 内容 \\
\midrule
系统名称 & 玉器门店顾客服务与经营管理系统 \\
目标组织 & 玉器销售门店柜台（档口老板 + 销售） \\
负责人 & 门店老板（柜台） \\
业务目标 & 提高门店销售额，降低货品周转周期 \\
度量指标 & ① 月度销售客户数；② 货品加权周转周期（统计口径：从货品进货入库开始，至货品销售完成并收回货款为止，按货品总金额加权平均） \\
对比时间 & 2025 年全年业务月均值（引入系统前，沿用纸质台账的做法） \\
改进值 & 销售额提升 10\%；货品周转周期下降 20\% \\
\bottomrule
\end{tabularx}
\end{table}

\begin{itemize}[leftmargin=1.6em, itemsep=2pt]
  \item 玉器一物一价、金额差异大，周转周期必须按件追踪：从进货入库开始，
        到销售完成并收回货款为止，按货品总金额加权平均。六位数的货压三四年是常事，
        几百到一两千的货一周多就能走，不加权的话平均值会掩盖高价货的真实积压。
        月度销售客户数按当月成交客户去重计。
  \item 改进值是拟议目标，不是承诺；周转周期下降 20\% 来自高价压柜货品及早发现和调剂，
        以系统记录为统计依据。销售额提升 10\% 主要来自老客户复购与线上引流到店的转化——
        收料成本和定价靠眼力，系统不决定成交价，只把货品与客户记录支撑起来。
\end{itemize}

\subsection{业务用例图}

业务用例图描述目标组织对外提供的价值。目标组织是玉器门店（档口），
业务执行者共七位：客户、回流出让方、玉料供应商、加工师傅、成品玉器供应商、鉴定机构、同行商户。
业务用例图见图 \ref{fig:01-business-usecase.png}；手绘稿见图 \ref{fig:01a-business-usecase-hand.jpg}。

\figab[0.86]{业务用例图（手绘稿）}{01a-business-usecase-hand.jpg}

\figab[0.9]{业务用例图（PlantUML 绘制）}{01-business-usecase.png}

引入目标系统后，新增了面向客户的线上服务用例（预定商品、预约看货）、面向合作商户的查看调货价目、面向供应商的提交出让申请与面向加工师傅的查看加工任务，改进业务用例图见图 \ref{fig:02-business-usecase-to-be.png}。

\figab[0.9]{改进业务用例图：新增的线上服务用例}{02-business-usecase-to-be.png}

\noindent
图中组织边界内共有 11 个业务用例，名称按业务执行者提出请求的口径命名：
面向客户侧的定制、购买玉器、退货，
面向货品来源侧的回收（收购回流玉器）、出售玉料、接加工、出售成品玉器、送检、调货，
以及承载货款结算的支取工费、支付供应商货款。
货品来源侧的四个方向（出售玉料、出售成品玉器、回收、调货）都对应同一条货款结算线，
所以支付供应商货款画成一个用例、由多方共同参与。

\subsection{现状业务序列图}

现状业务序列图描述组织当前的实际做法。
图中人员与记录工具都照实画（纸质销售台账、进货单据都是业务实体），
消息统一写成“A 请求 B 做某事”的服务调用，名字用“动词+宾语()”，
代表的是责任与目的，跟数据流向无关：箭头指向谁，就由谁承担这项工作，
完成后用虚线把结果返回给请求方，处理期间画激活条。
符号约定：业务执行者（客户）画小人；业务工人（销售、老板）画圆圈内人形，
圆圈右下加一道斜线切口、顶部带小箭头；业务实体（纸质销售台账、进货单据）画圆圈加右下斜线，
圆圈底下加一横；对象名统一写成加粗下划线的“:名字”。

现状的日常经营拆成场景逐一定格：现场议价成交、退货、成品进货、
同行调货（调入、调出各一张）、玉料采购与加工。
这些流程里，问题出在四处：
调货靠微信在图片上写价、发图成交，记录会丢、图片会被清理，事后认不出调出去的是哪件货，
对账只能靠回忆；售出只在台账记一笔金额，没有和具体货品关联，货品从入库到售出的时间无从对齐，
周转周期只能估算；客户成交后不留任何记录，月度销售客户数只能翻账本数；
月底汇总靠人工，忙时漏记当晚无从发现。

\figab[0.9]{现状场景：议价成交}{02-as-is-sale.png}
\figab[0.9]{现状场景：退货}{03-as-is-return.png}
\figab[0.9]{现状场景：成品进货}{04-as-is-goods-in.png}
\figab[0.9]{现状场景：同行调货（合作商户调货给我们）}{05-as-is-transfer-in.png}
\figab[0.9]{现状场景：同行调货（调出货品给合作商户）}{06-as-is-transfer-out.png}
\figab[0.9]{现状场景：玉料采购与加工}{07-as-is-material.png}

\subsection{改进业务序列图}

改进思路是把人工记账与统计交给新引入的系统：新系统是新的业务实体，
接手纸质销售台账与进货单据的部分责任。货品入库时逐件建档（编号、品名、来源、成本、入库时间），
成交时按货品编号登记并更新货品状态，销售与周转统计由系统按件自动累计；
保存、校验、统计都是系统内部行为，画成自消息，不单独画数据库生命线。

六个现状场景逐一改进：凡是改进场景里从外部指向目标系统的消息，
都映射为 \ref{fig:08-system-usecase.png} 的系统用例——登记成交、登记退货、登记货品入库、
预定商品、预约查看货品（图中消息为“预约看货”）、查看调货价格（消息为“查看调货价目”）、
提交货品出让申请、审核出让申请、预约上门时间、查看加工任务、提交完工交付申请，
以及由各货款消息聚合的登记货款记录、由调货往来与库存消息聚合的维护调货记录，共 13 个；
维护客户信息、登记送检与证书、查看销售与周转统计、维护基础数据、登录、浏览货品、
开通外部账户 7 个用例是支撑上述用例与愿景指标所必需的补充，在图中一并列出。

六个现状场景的改进版逐一如下：线上预定后到店成交、退货登记与货品状态回柜、
成品进货（供应商在线提交出让申请并预约上门）、同行调货（调入、调出）、
玉料采购与加工（加工师傅在线查看任务与提交完工交付）。
六个现状场景之外补一个新增业务场景——供应商寄售（改进场景图~\ref{fig:15-to-be-consignment.png}），
该场景原来不存在，没有现状图，直接给出改进版：
寄售货品凭合同建档入库，售出后按寄售价结算，到期未售出货品退回。

\figab[0.9]{改进场景：议价成交（线上预定 + 成交当场登记）}{09-to-be-sale.png}
\figab[0.9]{改进场景：退货}{10-to-be-return.png}
\figab[0.9]{改进场景：成品进货（在线出让申请与预约上门）}{11-to-be-goods-in.png}
\figab[0.9]{改进场景：同行调货（合作商户调货给我们）}{12-to-be-transfer-in.png}
\figab[0.9]{改进场景：同行调货（调出货品给合作商户）}{13-to-be-transfer-out.png}
\figab[0.9]{改进场景：玉料采购与加工（在线查看任务与交付）}{14-to-be-material.png}
\figab[0.9]{改进场景：供应商寄售（合同入库与寄售结算）}{15-to-be-consignment.png}

\subsection{业务改进对照}

\needspace{8\baselineskip}
\begin{table}[H]
\centering
\caption{现状痛点与改进后对照}
\begin{tabularx}{\textwidth}{@{}p{3.2cm} X X@{}}
\toprule
环节 & 现状 & 改进后 \\
\midrule
调货往来 & 微信发图写价，记录会丢，对账扯皮 & 调货往来与调货库存进系统，按货品编号与克重对账 \\
成交记录 & 手写小票 + 台账，只有金额 & 按货品编号登记成交，必填货品与成交价 \\
货品档案 & 无（进货信息靠单据和记忆） & 逐件建档：编号、品名、来源、成本、入库时间 \\
周转统计 & 无法统计 & 系统按件记录入库与售出回款，按金额加权计算周转周期 \\
客户与营销 & 散落在微信和记忆里 & 客户档案关联购买记录，支撑精准营销 \\
月度汇总 & 月底人工翻账本 & 自动汇总月度销售额与销售客户数 \\
\bottomrule
\end{tabularx}
\end{table}



% ============================================================
\section{B 需求}
% ============================================================

\subsection{系统用例}

用例（Use Case）即使用实例，从用户的角度出发描述计算机能够实现的具体目标。
系统用例图从改进业务序列图映射而来：系统用例 = 从外部指向目标系统的消息，
系统执行者 = 与目标系统实例有实线相连的实例所代表的系统，没有辅执行者。
其中登记成交、登记退货、登记货品入库、预定商品、预约查看货品、查看调货价格、
提交货品出让申请、审核出让申请、预约上门时间、查看加工任务、提交完工交付申请、
登记货款记录与维护调货记录 13 个用例直接来自六个改进业务场景的消息映射（见 A 部分改进业务序列图一节）；
登记寄售合同、结算寄售货款、结算调货台账、登记分期计划、记录分期还款与查看库存流水 6 个用例
来自新增的供应商寄售场景（图~\ref{fig:15-to-be-consignment.png}）与台账、分期、流水三项扩展；
维护客户信息、登记送检与证书、查看销售与周转统计、维护基础数据、登录、浏览货品、
开通外部账户 7 个用例是支撑上述用例与愿景指标所必需的补充。
系统的执行者分为老板、销售、客户、合作商户、供应商与加工师傅六类，
故按执行者端进行用例分析。总体用例图见图 \ref{fig:08-system-usecase.png}，
共 6 个执行者、26 个用例；
系统负责货品档案、成交记录、统计，面向客户的货品展示、预约与预定，
以及供应商与加工师傅的在线协同，不接入在线支付接口，不参与鉴定和物流。

\figab[0.60]{系统用例图：系统边界、执行者与用例}{08-system-usecase.png}

账号角色设老板、销售、合作商户、供应商与加工师傅五种，
合作商户、供应商与加工师傅的账号都由老板逐个开通（UC-12），不开放自助注册；
客户浏览、预约与预定无需账号，前端不展示价格——标出的价格并不是底价。
柜台端用电脑操作；供应商、加工师傅与同行商户多在散档口、不便用电脑且文化程度不一，
外部协同端做成微信小程序，操作步骤尽量少，账号与初始口令由老板线下交付。
入库与出库在现实里由老板和销售同时经手，改进后的口径是默认分工加系统兜底：
出库（成交、调出、退货）由销售登记，入库建档以老板为主，
老板不在店时由销售代理收货建档，系统保留经手人，记错能追到责任人；
送检由老板在系统提交货品清单、销售每天去鉴定机构集中办理，同样留经手人。
线上预约与上门预约统一排期：客户在系统约好到店时段后直接找对应销售，
避免销售当天已有接待任务时老板两头协调。
“登录”本身不产生业务价值，但除顾客端外，进入各端功能都要先验证身份，
单独列为用例（UC-08），由各端主用例以 \code{<<include>>} 指向。

按执行者端划分的用例如下。

\subsubsection{管理端用例（老板）}

管理端由老板使用。老板在系统中承担经营决策者的角色，拥有登录、登记货品入库、
查看销售与周转统计、维护基础数据、开通外部账户、审核出让申请、登记货款记录、
维护调货记录、登记寄售合同、结算寄售货款、结算调货台账与查看库存流水等用例，
具体用例图如图所示。

\figab[0.88]{管理端用例}{uc-boss.png}

用例详细说明如下：

\noindent（1）登录：老板使用由自己开通或设置的账号登录，验证通过后进入管理端功能。

\noindent（2）登记货品入库：对采购、回收、调货调入、加工完工与寄售的货品逐件建档，
录入来源、成本与最低售价，货品状态置为“在柜”，周转周期从这里开始计时。

\noindent（3）查看销售与周转统计：系统按成交记录自动累计月度销售额、销售客户数
与货品加权周转周期，老板随时查看，是愿景指标的读数入口。

\noindent（4）维护基础数据：维护来源方（供应商/加工师傅/同行商户）、鉴定机构与货品类别，
支撑其余用例的参照数据。

\noindent（5）开通外部账户：为合作商户、供应商与加工师傅逐个开通账号与权限，
停用即收回，外部角色不能自行注册。

\noindent（6）审核出让申请：审核供应商在线提交的货品出让申请，
通过的转入入库流程，驳回的填写原因并留痕。

\noindent（7）登记货款记录：登记采购货款、调货货款、加工工费与寄售货款的新增与收付，
每笔货款可追溯到来源单据与经手人。

\noindent（8）维护调货记录：维护调货往来记录与调货库存记录，
调货货品去向可追溯，结算不扯皮。

\noindent（9）登记寄售合同：供应商寄存货品凭合同建档入库，
寄售价记为成本作为结算依据，合同与明细一次登记。

\noindent（10）结算寄售货款：已售出的寄售货品按寄售价与供应商结算，
明细置为已结算，回款有结算单为据。

\noindent（11）结算调货台账：按商户汇总未结算调货记录一次对账，
结算后单据置为已结算，鉴赏期届满即结清。

\noindent（12）查看库存流水：任一货品的每次状态变更都留有流水，
进出有账可查，是调货库存一本账的追溯底账。

\subsubsection{柜台端用例（销售）}

柜台端由销售使用。销售在系统中承担柜台经手的角色，拥有登录、登记成交、
登记退货、维护客户信息、登记送检与证书、登记分期计划与记录分期还款等用例，
具体用例图如图所示。

\figab[0.88]{柜台端用例}{uc-sale.png}

用例详细说明如下：

\noindent（1）登录：销售使用老板开通的账号登录，验证通过后进入柜台端功能。

\noindent（2）登记成交：成交（全款或定金）当场登记到货品编号上，
系统校验货品在柜并更新货品状态，销售额与销售客户数两个度量都从这里来。

\noindent（3）登记退货：收回退货，货品状态改回“在柜”，冲减销售统计，
保证销售与库存数据准确。

\noindent（4）维护客户信息：维护客户姓名、联系方式、偏好与购买记录，
是精准营销的前提。

\noindent（5）登记送检与证书：登记送检记录与鉴定证书信息并挂到对应货品，
客户索证时随时可调，是高价货品销售的依据。

\noindent（6）登记分期计划：大额成交当场谈定分期，系统按成交价生成分期计划，
总额与期数一次确定。

\noindent（7）记录分期还款：录入每期还款并推进到期日，
结清自动归档，逾期计划随时可查。

\subsubsection{顾客端用例（客户）}

顾客端面向到店与线上客户。客户无需登录即可使用，拥有浏览货品、预约查看货品、
预定商品等用例，具体用例图如图所示。

\figab[0.88]{顾客端用例}{uc-cust.png}

用例详细说明如下：

\noindent（1）浏览货品：无需登录浏览在柜货品的图文展示，不显示价格，
把在柜货品摆到前端引流。

\noindent（2）预约查看货品：提交联系方式与到店时段，预约到店看货，转化到店。

\noindent（3）预定商品：留手机号预定在柜货品，一个手机号 3 次额度、
如约达成后恢复，锁定到店意向。

\subsubsection{商户端用例（合作商户）}

商户端由合作商户使用。合作商户拥有登录、查看调货价格等用例，具体用例图如图所示。

\figab[0.88]{商户端用例}{uc-merch.png}

用例详细说明如下：

\noindent（1）登录：合作商户使用老板开通的账号登录，验证通过后仅可见调货价目。

\noindent（2）查看调货价格：登录后查看在柜货品的调货价目，
减少问价与开户的人工往返。

\subsubsection{供应端用例（供应商）}

供应端由供应商使用（含玉料供应商、成品供应商与回流出让方）。供应商拥有登录、
提交货品出让申请、预约上门时间等用例，具体用例图如图所示。

\figab[0.88]{供应端用例}{uc-supp.png}

用例详细说明如下：

\noindent（1）登录：供应商使用老板开通的账号登录，验证通过后进入供应端功能。

\noindent（2）提交货品出让申请：在线提交玉料、成品或回流货品的出让申请，
货源线索集中在线上，不用上门先谈一轮。

\noindent（3）预约上门时间：申请通过后预约上门交接时段，按约前来，不白跑。

\subsubsection{加工端用例（加工师傅）}

加工端由加工师傅使用。加工师傅拥有登录、查看加工任务、提交完工交付申请、预约上门时间等用例，
具体用例图如图所示。

\figab[0.88]{加工端用例}{uc-work.png}

用例详细说明如下：

\noindent（1）登录：加工师傅使用老板开通的账号登录，验证通过后进入加工端功能。

\noindent（2）查看加工任务：登录查看分派给自己的加工任务、加工要求与工费，
不用电话逐个通知。

\noindent（3）提交完工交付申请：完工后在线报备，等待老板确认并安排交接入库，
完工信息第一时间可见。

\noindent（4）预约上门时间：在线预约上门交接时段，出让交货与完工交付都不用电话约定。

\subsection{用例清单}

\needspace{8\baselineskip}
\begin{table}[H]
\centering
\small
\caption{系统用例清单（排序依据是愿景）}
\begin{tabularx}{\textwidth}{@{}l l X c@{}}
\toprule
编号 & 用例 & 说明 & 排序 \\
\midrule
UC-01 & 登记成交 & 成交（全款或定金）当场登记到货品编号上，系统校验并保存 & 1（关键） \\
UC-02 & 登记货品入库 & 采购、回收、调货、加工完工、寄售的货品逐件建档入库 & 2（关键） \\
UC-03 & 登记退货 & 收回退货，货品状态改回“在柜”，冲减销售统计 & 4 \\
UC-04 & 维护客户信息 & 客户姓名、联系方式、偏好与购买记录，支撑精准营销 & 5 \\
UC-05 & 登记送检与证书 & 送检记录与鉴定证书信息挂到对应货品 & 6 \\
UC-06 & 查看销售与周转统计 & 月度销售额、销售客户数、货品加权周转周期 & 3 \\
UC-07 & 维护基础数据 & 来源方（供应商/加工师傅/同行商户）、鉴定机构与货品类别 & 7 \\
UC-08 & 登录 & 验证用户名与口令后按角色进入系统（被 UC-01～UC-07 与 UC-11、UC-12、UC-14～UC-26 包含） & — \\
UC-09 & 浏览货品 & 客户无需登录浏览在柜货品图文展示，不显示价格 & 8 \\
UC-10 & 预约查看货品 & 客户提交联系方式与到店时段，预约到店看货 & 9 \\
UC-11 & 查看调货价格 & 合作商户登录后查看在柜货品的调货价目 & 10 \\
UC-12 & 开通外部账户 & 老板为合作商户、供应商与加工师傅逐个开通账号与权限 & 11 \\
UC-13 & 预定商品 & 客户留手机号预定在柜货品，一个手机号 3 次额度、如约达成恢复 & 12 \\
UC-14 & 提交货品出让申请 & 供应商（含回流出让方）在线提交货品出让申请 & 13 \\
UC-15 & 预约上门时间 & 供应商与加工师傅在线预约上门交接时段 & 14 \\
UC-16 & 查看加工任务 & 加工师傅登录查看分派给自己的加工任务（含要求、工费与期限） & 15 \\
UC-17 & 提交完工交付申请 & 加工师傅完工后在线提交交付申请，等待老板确认并安排入库 & 16 \\
UC-18 & 审核出让申请 & 老板审核出让申请，通过后转入入库流程 & 17 \\
UC-19 & 登记货款记录 & 供货方货款、调货货款与加工工费的新增与收付记录 & 18 \\
UC-20 & 维护调货记录 & 调货往来记录与调货库存记录的新增与更新 & 19 \\
UC-21 & 登记寄售合同 & 供应商寄存货品逐件建档入库，寄售价作结算依据 & 20 \\
UC-22 & 结算寄售货款 & 已售寄售货品按寄售价与供应商结算，明细置已结算 & 21 \\
UC-23 & 结算调货台账 & 按商户汇总未结算调货记录，一次对账结清 & 22 \\
UC-24 & 登记分期计划 & 大额成交分期付款，总额等于成交价 & 23 \\
UC-25 & 记录分期还款 & 录入还款、推进到期日，结清自动归档 & 24 \\
UC-26 & 查看库存流水 & 货品每次状态变更留痕，进出有账可查 & 25 \\
\bottomrule
\end{tabularx}
\end{table}

\noindent
逐项理由如下：UC-01 登记成交排第 1，销售额、销售客户数两个度量都从成交记录来；
UC-02 登记货品入库排第 2，周转周期从入库开始计时，没有入库记录的货品无法追踪；
UC-06 查看销售与周转统计排第 3，愿景指标的读数入口；
UC-03 退货排第 4，直接影响销售与库存数据的准确性；
UC-04 维护客户信息排第 5，精准营销的前提；UC-05 送检与证书排第 6，高价货品销售的依据；
UC-07 基础数据排第 7，支撑前述用例的参照数据；UC-08 被其他用例包含，不单独排序。
UC-09、UC-10 把在柜货品摆到前端引流并转化到店，排在台账内功之后；
UC-11、UC-12 服务同行调货，减少问价与开户的人工往返；
UC-13 预定商品锁定到店意向，UC-14～UC-18 把供应商、回流出让方与加工师傅的往来搬到线上，
省去电话问货与跑腿上门的时间；
UC-19、UC-20 承接改进场景里的记账消息，货款与调货一本账；
UC-21～UC-23 把寄售与台账结算搬上线，寄售价作结算依据、按商户一次对账，
钱账两清有据可查；UC-24～UC-25 承接大额成交的分期付款，逾期可追踪；
UC-26 库存流水给一本账补上追溯底账，货品每次状态变更留痕。

\subsection{系统用例规约}

根据各用例设计用例规约，进一步详细规定用例的用例编号与用例名、执行者、前置条件、
后置条件、涉众利益、基本路径、扩展，以及字段列表、业务规则、质量需求、设计约束
四个补充约束，每个规约用一张两列表格，下文按 UC-01～UC-26 逐个给出。
前置与后置条件只写系统能检测到的约束，扩展路径只写系统能感知并要处理的意外与分支。

\noindent（1）登记成交

{\small
\begin{longtable}{@{}p{2.4cm} p{12.9cm}@{}}
\caption{登记成交用例规约}\label{tab:uc-01}\\
\toprule
\textbf{标题} & \textbf{内容} \\
\midrule
\endfirsthead
\toprule
\textbf{标题} & \textbf{内容} \\
\midrule
\endhead
用例编号：用例名 & UC-01：登记成交 \\
\midrule
执行者 & 销售（主执行者） \\
\midrule
前置条件 & 系统能检测：所选货品已建档入库且处于“在柜”状态（散客可不关联客户档案） \\
\midrule
后置条件 & 新增一条成交记录，必然关联到唯一货品与经手销售；货品状态改为“已售”或“已订” \\
\midrule
涉众利益 & 销售：登记要快、不打断成交；\newline 老板：每笔成交可追溯到货品、客户与时间，周转与营销数据可信；\newline 客户：价格与小票一致，售后有据 \\
\midrule
基本路径 & 1. 销售扫描货品条形码或提供货品编号、成交价、支付方式（全款/定金）与客户（散客可不提供）；\newline 2. 系统校验货品处于“在柜”状态；\newline 3. 系统校验成交价合法（不低于该货品最低售价）；\newline 4. 系统创建成交记录，货品状态改为“已售”（定金成交则标记“已订”）；\newline 5. 系统反馈登记结果 \\
\midrule
扩展 & 2a. 货品编号不存在或状态非“在柜”（已订、已售、预留、已调出、代销中、已送检、待处理）：2a1 系统反馈应重新选择货品；2a2 返回1；\newline 3a. 成交价低于最低售价：3a1 系统要求录入老板授权价及授权说明，无授权时反馈应重新报价；3a2 返回1；\newline 3b. 定金金额超过成交价：3b1 系统反馈应重新输入；3b2 返回1；\newline 4a. 保存失败：4a1 系统反馈应稍后重试；4a2 返回1 \\
\midrule
字段列表 & 货品 = 货品编号 + 品名 + 类别 + 料种（自定义） + 产地（自定义） + 器型（自定义） + 产状{山料|籽料|山流水料|戈壁料} + 克重 + 尺寸 + 证书号 + 图片 + 来源类型 + 来源方 + 成本 + 最低售价 + 入库时间 + 状态{在柜|已送检|预留|已售|已订|已调出|代销中|待处理}；\newline 成交记录 = 成交编号 + 成交时间 + 货品 + 成交价 + 支付方式{全款|定金} + 已收金额 + 客户 + 经手人 \\
\midrule
业务规则 & ①每件玉器一物一码，一件货品同时只能有一条有效成交记录；\newline ②成交价不得低于该货品最低售价，低于时必须录入老板授权；\newline ③客户明确要求预留但尚未付款的，货品状态置为“预留”，预留期间不得出售给他人；\newline ④客户选择分期留货的，收定金即标记“已订”，货品留店不出库，与预留同样不得再售他人，尾款付清并交付出库后才计为完成销售；\newline ⑤寄卖在同行处的货品状态为“代销中”，货权仍归本店，售出后转为“已售”并与同行结算；\newline ⑥成交时间以系统时间为准，发生时间与录入时间分别记录 \\
\midrule
质量需求 & 可用性：完成一次登记不超过 4 次操作、10 秒内完成；\newline 性能：提交后 1 秒内反馈 \\
\midrule
设计约束 & 柜台端用电脑操作并部署在门店本地网络；供应商、加工师傅与同行商户通过微信小程序访问；货款当面收付（现金或个人收款码），可按需要配小票打印机与条形码标签打印机；不接入在线支付接口 \\
\midrule
\bottomrule
\end{longtable}}
\noindent（2）登记货品入库

{\small
\begin{longtable}{@{}p{2.4cm} p{12.9cm}@{}}
\caption{登记货品入库用例规约}\label{tab:uc-02}\\
\toprule
\textbf{标题} & \textbf{内容} \\
\midrule
\endfirsthead
\toprule
\textbf{标题} & \textbf{内容} \\
\midrule
\endhead
用例编号：用例名 & UC-02：登记货品入库 \\
\midrule
执行者 & 老板（主执行者） \\
\midrule
前置条件 & 来源方（供应商、加工师傅、同行商户）已在基础数据中维护 \\
\midrule
后置条件 & 新增一条货品档案，状态“在柜”，入库时间以系统时间为准 \\
\midrule
涉众利益 & 老板：每件货品都有唯一编号和完整来源、成本信息，周转统计才有依据；\newline 销售：录入成交时按编号即可带出货品资料 \\
\midrule
基本路径 & 1. 老板提供货品信息（品名、类别、料种、产地、器型、克重、尺寸、证书号、图片、来源类型、来源方、成本、最低售价）；\newline 2. 系统校验字段与金额合法；\newline 3. 系统逐件生成货品编号、保存档案并打印条形码标签（印货品编号，系挂于货品）；\newline 4. 系统反馈入库结果与编号清单 \\
\midrule
扩展 & 1a. 来源类型为回收：1a1 系统按回流作价确定成本并标记来源与作价；\newline 1b. 来源类型为调货：1b1 系统标记来源与关联商户，不记成本、只记往来；\newline 2a. 必填字段缺失或金额不合法：2a1 系统反馈应重新输入；2a2 返回1；\newline 3a. 同批次多件（如珠串类）：3a1 系统逐件建档、共用批次信息，编号各不相同 \\
\midrule
字段列表 & 货品来源 = 来源类型{采购|回收|调货|加工完工|寄售} + 来源方 + 来源单据；\newline 入库记录 = 入库时间 + 验收人 + 货品清单 \\
\midrule
业务规则 & ①一件货品一个编号，重件需拆分为多条档案；\newline ②回收货品按回流作价入成本；\newline ③调货调入不记成本，只记往来；\newline ④加工完工按料加工合计入成本；\newline ⑤寄售货品按寄售价入成本，作为结算依据 \\
\midrule
质量需求 & 可用性：逐件建档不超过 30 秒/件 \\
\midrule
设计约束 & 与现有纸质台账并行运行一个月，双轨核对无误后停用台账 \\
\midrule
\bottomrule
\end{longtable}}
\noindent（3）登记退货

{\small
\begin{longtable}{@{}p{2.4cm} p{12.9cm}@{}}
\caption{登记退货用例规约}\label{tab:uc-03}\\
\toprule
\textbf{标题} & \textbf{内容} \\
\midrule
\endfirsthead
\toprule
\textbf{标题} & \textbf{内容} \\
\midrule
\endhead
用例编号：用例名 & UC-03：登记退货 \\
\midrule
执行者 & 销售（主执行者） \\
\midrule
前置条件 & 待退货的成交记录存在且未标记过退货 \\
\midrule
后置条件 & 原成交记录标记“已退货”，货品状态改回“在柜”，当期销售统计冲减 \\
\midrule
涉众利益 & 客户：真假与质量问题能退、退款及时，不喜欢可协商换货；\newline 老板：退货限于成交后 7 天内且限真假与质量原因，防止随意退货冲减业绩；\newline 销售：不用手工调整台账，换货协商有记录可依 \\
\midrule
基本路径 & 1. 销售提供原成交记录与退货原因；\newline 2. 系统校验该成交存在且未标记过退货，且在成交后 7 天内、退货原因属于可退范围；\newline 3. 系统新增退货记录，货品状态改回“在柜”；\newline 4. 系统反馈登记结果 \\
\midrule
扩展 & 2a. 原成交记录不存在或已退货：2a1 系统反馈不允许重复退货；\newline 2b. 已超过成交后 7 天：2b1 系统反馈已超过退货期限，不予办理；\newline 2c. 退货原因不属于可退范围：2c1 系统反馈该原因不在可退范围内，不予办理；\newline 3a. 退货价与原成交价不一致：3a1 销售说明差额原因，系统将差额计入退货记录 \\
\midrule
字段列表 & 退货记录 = 退货时间 + 原成交记录 + 退货原因{真假存疑|质量问题} + 退款金额 + 经手人 \\
\midrule
业务规则 & ①退货必须关联原成交记录，无记录不走系统退货；\newline ②退货期限自成交日起 7 天，超期不予办理；\newline ③仅真假存疑或质量问题可退；圈口不合、主观不满意等原因不予退款，由销售协商换货处理；\newline ④退货后货品状态改回“在柜”，可再次出售；\newline ⑤当月退货冲减当月销售统计 \\
\midrule
质量需求 & 可用性：完成一次退货登记不超过 3 次操作 \\
\midrule
设计约束 & 退款为当面现金退回或按原收款方式退回，不接入在线支付接口 \\
\midrule
\bottomrule
\end{longtable}}
\noindent（4）维护客户信息

{\small
\begin{longtable}{@{}p{2.4cm} p{12.9cm}@{}}
\caption{维护客户信息用例规约}\label{tab:uc-04}\\
\toprule
\textbf{标题} & \textbf{内容} \\
\midrule
\endfirsthead
\toprule
\textbf{标题} & \textbf{内容} \\
\midrule
\endhead
用例编号：用例名 & UC-04：维护客户信息 \\
\midrule
执行者 & 销售（主执行者） \\
\midrule
前置条件 & 新增客户时可提供姓名与联系方式 \\
\midrule
后置条件 & 客户档案新增或更新，客户编号唯一；历史成交记录中的客户引用保持不变 \\
\midrule
涉众利益 & 销售：回头客到店能快速调出资料与购买记录，手机号可一键复制去加微信；\newline 老板：客户资料完整真实，来源渠道与标签能看出哪个渠道有效、谁是核心客户；\newline 客户：个人信息只用于门店经营，不外泄 \\
\midrule
基本路径 & 1. 销售提供客户姓名与联系方式，并选择来源渠道与标签；\newline 2. 系统校验手机号未与其他客户重复；\newline 3. 系统保存客户档案并生成客户编号；\newline 4. 系统反馈维护结果 \\
\midrule
扩展 & 2a. 手机号已存在：2a1 系统带出原客户档案，销售可补充修改；2a2 返回1；\newline 3a. 必填项缺失：3a1 系统反馈应补全后保存；3a2 返回1 \\
\midrule
字段列表 & 客户 = 客户编号 + 姓名 + 手机号 + 微信号 + 来源渠道{闲鱼|微信朋友圈|朋友介绍|线下|直播|老客户复购|其他} + 标签{大客户|懂行|爽快|要还价} + 偏好（品类+圈口+预算）+ 购买记录 + 跟进记录 + 备注 \\
\midrule
业务规则 & ①同一手机号视为同一客户，不重复建档；\newline ②购买记录由成交登记自动带出，不手工录入；\newline ③来源渠道与标签从固定选项中选择，不自由填写，以便按渠道统计 \\
\midrule
质量需求 & 可用性：按编号、姓名、手机号、微信号、标签任一条件查询，3 秒内反馈；\newline 安全性：手机号等联系方式仅老板与销售可见 \\
\midrule
设计约束 & 不做短信/微信营销推送，仅作查询与筛选；不做会员等级、积分与优惠券；不做客户画像算法打分，核心客户按累计成交与标签筛选得出 \\
\midrule
\bottomrule
\end{longtable}}
\noindent（5）登记送检与证书

{\small
\begin{longtable}{@{}p{2.4cm} p{12.9cm}@{}}
\caption{登记送检与证书用例规约}\label{tab:uc-05}\\
\toprule
\textbf{标题} & \textbf{内容} \\
\midrule
\endfirsthead
\toprule
\textbf{标题} & \textbf{内容} \\
\midrule
\endhead
用例编号：用例名 & UC-05：登记送检与证书 \\
\midrule
执行者 & 销售（主执行者） \\
\midrule
前置条件 & 货品已建档入库且状态为“在柜”或“预留” \\
\midrule
后置条件 & 新增一条送检记录，货品状态改为“已送检”；证书信息保存后与货品关联，货品状态恢复送检前的状态（在柜或预留） \\
\midrule
涉众利益 & 销售：证书编号与货品一一对应，客户索证时随时可调，销售时直接出示证书；\newline 老板：不用亲自跑送检，只需安排与把关；\newline 鉴定机构：送检信息与机构回执一致 \\
\midrule
基本路径 & 1. 销售提供货品编号、送检机构与送检日期；\newline 2. 系统校验货品状态为“在柜”或“预留”；\newline 3. 系统新增送检记录，货品状态改为“已送检”；\newline 4. 销售录入鉴定机构返回的证书编号与检测结论；\newline 5. 系统保存证书信息并关联货品，货品状态恢复送检前状态；\newline 6. 系统反馈登记结果 \\
\midrule
扩展 & 2a. 货品状态非“在柜”且非“预留”（已售、已调出、代销中）：2a1 系统反馈应重新选择货品；2a2 返回1；\newline 4a. 检测结论为不合格：4a1 系统保存结论并把货品状态改为“待处理”，由老板决定退回供应商或停售 \\
\midrule
字段列表 & 送检记录 = 送检时间 + 货品 + 送检机构 + 送检人；\newline 证书 = 证书编号 + 鉴定机构 + 检测结论 + 签发日期 + 证书二维码（机构印制，扫码可查线上鉴定证书） \\
\midrule
业务规则 & ①证书编号与货品一一对应，不得复用；\newline ②送检期间货品状态为“已送检”，不得出售；\newline ③证书编号录入后不可修改，只能作废重录；\newline ④鉴定证书只判真假，产地由老板凭眼力判断并记录在货品档案 \\
\midrule
质量需求 & 性能：提交后 1 秒内反馈 \\
\midrule
设计约束 & 鉴定机构信息从基础数据带出，不与机构系统对接 \\
\midrule
\bottomrule
\end{longtable}}
\noindent（6）查看销售与周转统计

{\small
\begin{longtable}{@{}p{2.4cm} p{12.9cm}@{}}
\caption{查看销售与周转统计用例规约}\label{tab:uc-06}\\
\toprule
\textbf{标题} & \textbf{内容} \\
\midrule
\endfirsthead
\toprule
\textbf{标题} & \textbf{内容} \\
\midrule
\endhead
用例编号：用例名 & UC-06：查看销售与周转统计 \\
\midrule
执行者 & 老板（主执行者） \\
\midrule
前置条件 & 系统内已有货品入库记录与成交记录 \\
\midrule
后置条件 & 只读查询，不改变任何数据 \\
\midrule
涉众利益 & 老板：统计数字能与纸质台账核对一致，按月对比销售额与周转周期两个愿景指标；\newline 销售：可查自己经手的成交，数据不被他人改动 \\
\midrule
基本路径 & 1. 老板选择统计时段（月/季/年）；\newline 2. 系统汇总该时段销售额、销售客户数与货品加权周转周期；\newline 3. 系统按品种与价位筛选周转最慢的货品；\newline 4. 系统按来源渠道汇总成交，并列出累计贡献靠前的客户；\newline 5. 系统反馈统计结果 \\
\midrule
扩展 & 2a. 时段内无记录：2a1 系统反馈该时段无数据；2a2 返回1 \\
\midrule
字段列表 & 统计口径 = 月度销售额 + 销售客户数（去重）+ 货品加权周转周期（按货品总金额加权）+ 来源渠道成交分布 + 核心客户贡献（累计成交额与占比） \\
\midrule
业务规则 & ①销售额按成交时间归属月份，退货冲减；\newline ②周转周期从货品入库到销售完成并收回货款，按货品总金额加权 \\
\midrule
质量需求 & 性能：统计查询 2 秒内反馈 \\
\midrule
设计约束 & 统计界面只读，不提供数据导出 \\
\midrule
\bottomrule
\end{longtable}}
\noindent（7）维护基础数据

{\small
\begin{longtable}{@{}p{2.4cm} p{12.9cm}@{}}
\caption{维护基础数据用例规约}\label{tab:uc-07}\\
\toprule
\textbf{标题} & \textbf{内容} \\
\midrule
\endfirsthead
\toprule
\textbf{标题} & \textbf{内容} \\
\midrule
\endhead
用例编号：用例名 & UC-07：维护基础数据 \\
\midrule
执行者 & 老板（主执行者） \\
\midrule
前置条件 & 无（身份验证由 UC-08 提供） \\
\midrule
后置条件 & 来源方（供应商/加工师傅/同行商户）、鉴定机构与货品类别等参照数据新增或更新，供入库与成交登记选择 \\
\midrule
涉众利益 & 老板：参照数据齐全，录入时直接选择、不用每次手输；\newline 销售：选择项稳定统一，统计不因写法不同出现偏差 \\
\midrule
基本路径 & 1. 老板选择要维护的数据类型（来源方/鉴定机构/货品类别，来源方以类别区分供应商/加工师傅/同行商户）；\newline 2. 老板提供新增、修改或删除的内容；\newline 3. 系统校验名称在同一类别内未重复；\newline 4. 系统保存并反馈维护结果 \\
\midrule
扩展 & 3a. 名称重复：3a1 系统反馈应更换名称；3a2 返回1；\newline 4a. 删除已被引用的数据项：4a1 系统反馈该数据项已被成交或入库记录引用，只能停用；4a2 返回1 \\
\midrule
字段列表 & 基础数据 = 来源方（类别 + 名称 + 联系人 + 联系电话） + 鉴定机构 + 货品类别 \\
\midrule
业务规则 & ①基础数据被引用后不可物理删除，只能停用；\newline ②供应商类型区分玉料/成品/回流 \\
\midrule
质量需求 & 可用性：新增一条基础数据不超过 3 次操作 \\
\midrule
设计约束 & 基础数据仅老板可维护，销售只读 \\
\midrule
\bottomrule
\end{longtable}}
\noindent（8）登录

{\small
\begin{longtable}{@{}p{2.4cm} p{12.9cm}@{}}
\caption{登录用例规约}\label{tab:uc-08}\\
\toprule
\textbf{标题} & \textbf{内容} \\
\midrule
\endfirsthead
\toprule
\textbf{标题} & \textbf{内容} \\
\midrule
\endhead
用例编号：用例名 & UC-08：登录 \\
\midrule
执行者 & 老板、销售、合作商户、供应商、加工师傅（被 UC-01～UC-07 与 UC-11、UC-12、UC-14～UC-26 以 \code{<<include>>} 包含的公共用例；客户无需登录） \\
\midrule
前置条件 & 用户账号已由老板建立 \\
\midrule
后置条件 & 登录成功后系统按角色开放对应功能；后续关键操作归属到该用户名下 \\
\midrule
涉众利益 & 老板：职责分离，销售无法使用授权价、基础数据维护等老板功能；\newline 销售：口令自行保管，经手记录归属清晰 \\
\midrule
基本路径 & 1. 执行者提供用户名与口令；\newline 2. 系统校验用户名存在且口令正确；\newline 3. 系统按角色开放对应功能并反馈登录结果 \\
\midrule
扩展 & 2a. 用户名不存在或口令错误：2a1 系统反馈应重试；2a2 返回1；\newline 2b. 同一账号连续 5 次口令错误：2b1 系统锁定账号 10 分钟，并提示联系老板解锁 \\
\midrule
字段列表 & 用户 = 用户名 + 姓名 + 口令（加密保存）+ 角色{老板|销售|合作商户|供应商|加工师傅} + 状态{启用|停用} \\
\midrule
业务规则 & ①连续 5 次口令错误锁定 10 分钟；\newline ②角色决定登录后的功能范围 \\
\midrule
质量需求 & 安全性：口令加密保存；\newline 性能：登录验证 1 秒内反馈 \\
\midrule
设计约束 & 不做自助注册，账号全部由老板开通 \\
\midrule
\bottomrule
\end{longtable}}
\noindent（9）浏览货品

{\small
\begin{longtable}{@{}p{2.4cm} p{12.9cm}@{}}
\caption{浏览货品用例规约}\label{tab:uc-09}\\
\toprule
\textbf{标题} & \textbf{内容} \\
\midrule
\endfirsthead
\toprule
\textbf{标题} & \textbf{内容} \\
\midrule
\endhead
用例编号：用例名 & UC-09：浏览货品 \\
\midrule
执行者 & 客户（无需登录，也不提供注册） \\
\midrule
前置条件 & 系统内存在状态为“在柜”且允许展示的货品档案 \\
\midrule
后置条件 & 只读查询，不改变任何数据 \\
\midrule
涉众利益 & 客户：随时了解门店在售货品，到店前有预期；\newline 老板：展示范围可控，价格信息不外露；\newline 销售：预约到店的客户带着目标来，接待更省力 \\
\midrule
基本路径 & 1. 客户请求浏览在柜货品；\newline 2. 系统检索状态为“在柜”且允许展示的货品图文信息；\newline 3. 客户选择货品查看详情；\newline 4. 系统反馈货品展示信息（不含任何价格） \\
\midrule
扩展 & 2a. 货品已售出、已调出或已送检：2a1 系统不再展示该货品 \\
\midrule
字段列表 & 展示项 = 图片 + 品名 + 类别 + 尺寸/圈口 + 描述（不含价格） \\
\midrule
业务规则 & ①仅展示“在柜”货品；\newline ②一物一码对应一页展示 \\
\midrule
质量需求 & 性能：图片加载 2 秒内完成；\newline 可用性：微信小程序或手机浏览器可直接打开 \\
\midrule
设计约束 & 前端不显示价格与成本 \\
\midrule
\bottomrule
\end{longtable}}
\noindent（10）预约查看货品

{\small
\begin{longtable}{@{}p{2.4cm} p{12.9cm}@{}}
\caption{预约查看货品用例规约}\label{tab:uc-10}\\
\toprule
\textbf{标题} & \textbf{内容} \\
\midrule
\endfirsthead
\toprule
\textbf{标题} & \textbf{内容} \\
\midrule
\endhead
用例编号：用例名 & UC-10：预约查看货品 \\
\midrule
执行者 & 客户（无需登录） \\
\midrule
前置条件 & 目标货品处于“在柜”且在展示范围内 \\
\midrule
后置条件 & 新增一条预约记录，关联货品、联系方式与期望到店时段；货品状态不变 \\
\midrule
涉众利益 & 客户：到店前确认货品在店、有人接待；\newline 销售：提前知道接待对象与看点；\newline 老板：预约数量是前端引流效果的直接指标 \\
\midrule
基本路径 & 1. 客户选择货品并请求预约查看；\newline 2. 客户提供姓名、手机号与期望到店时段；\newline 3. 系统校验手机号与时段合法，保存预约记录并生成预约编号；\newline 4. 系统反馈预约结果 \\
\midrule
扩展 & 2a. 手机号格式不合法：2a1 系统反馈应重新输入；2a2 返回2；\newline 3a. 同一手机号当日重复预约同一货品：3a1 系统提示已有预约，可查看原预约 \\
\midrule
字段列表 & 预约 = 预约编号 + 货品 + 姓名 + 手机号 + 期望时段 + 状态 \\
\midrule
业务规则 & ①同一货品同一时段允许多组预约；\newline ②预约信息仅用于到店接待 \\
\midrule
质量需求 & 性能：提交后 1 秒内反馈 \\
\midrule
设计约束 & 无需登录，手机号作为联系标识 \\
\midrule
\bottomrule
\end{longtable}}
\noindent（11）查看调货价格

{\small
\begin{longtable}{@{}p{2.4cm} p{12.9cm}@{}}
\caption{查看调货价格用例规约}\label{tab:uc-11}\\
\toprule
\textbf{标题} & \textbf{内容} \\
\midrule
\endfirsthead
\toprule
\textbf{标题} & \textbf{内容} \\
\midrule
\endhead
用例编号：用例名 & UC-11：查看调货价格 \\
\midrule
执行者 & 合作商户（主执行者，账号由老板逐个开通） \\
\midrule
前置条件 & 无（身份验证由 UC-08 提供） \\
\midrule
后置条件 & 只读查询，不改变任何数据 \\
\midrule
涉众利益 & 合作商户：自助掌握可调货品与调货价格，减少往返问价；\newline 老板：调货价格只对开通账户的同行可见，价格体系可控；\newline 客户：不受影响，前端不展示价格 \\
\midrule
基本路径 & 1. 合作商户选择货品类别或输入品名关键字；\newline 2. 系统检索状态为“在柜”的货品与对应调货价格；\newline 3. 系统反馈查询结果 \\
\midrule
扩展 & 2a. 货品未开放展示或已调出：2a1 系统不反馈该货品 \\
\midrule
字段列表 & 调货价目 = 货品 + 类别 + 圈口/尺寸 + 调货价 \\
\midrule
业务规则 & ①调货价仅对“合作商户”角色可见；\newline ②调货价与零售价分离维护 \\
\midrule
质量需求 & 性能：查询 2 秒内反馈 \\
\midrule
设计约束 & 价目页只读，不提供下载 \\
\midrule
\bottomrule
\end{longtable}}
\noindent（12）开通外部账户

{\small
\begin{longtable}{@{}p{2.4cm} p{12.9cm}@{}}
\caption{开通外部账户用例规约}\label{tab:uc-12}\\
\toprule
\textbf{标题} & \textbf{内容} \\
\midrule
\endfirsthead
\toprule
\textbf{标题} & \textbf{内容} \\
\midrule
\endhead
用例编号：用例名 & UC-12：开通外部账户 \\
\midrule
执行者 & 老板（主执行者） \\
\midrule
前置条件 & 开通对象已在基础数据中维护（来源方以类别区分供应商、加工师傅等） \\
\midrule
后置条件 & 新增一个外部账号并分配角色（可多选），可登录办理各自事项；停用后无法登录 \\
\midrule
涉众利益 & 老板：给谁开账号自己决定，价格与往来信息可控；\newline 合作商户：拿到账号即可自助查价，不用每次电话问价；\newline 供应商与加工师傅：线上提交申请与预约，减少往返沟通 \\
\midrule
基本路径 & 1. 老板选择要开通的角色（可多选）并提供账号信息（用户名与初始口令）；\newline 2. 系统校验用户名未重复；\newline 3. 系统创建账号并分配所选角色；\newline 4. 系统反馈开通结果 \\
\midrule
扩展 & 2a. 用户名已存在：2a1 系统反馈应更换用户名；2a2 返回1 \\
\midrule
字段列表 & 用户 = 用户名 + 姓名 + 口令（加密保存）+ 状态{启用|停用}；\newline 用户-角色 = 用户 + 角色{合作商户|供应商|加工师傅等，可多选} \\
\midrule
业务规则 & ①外部账号由老板逐个开通与停用，不能自行注册；\newline ②合作商户角色只有查看调货价格的权限，供应商与加工师傅角色只有各自协同事项的权限，均无内部管理功能；\newline ③口令加密保存，首次登录可修改 \\
\midrule
质量需求 & 安全性：口令加密保存；\newline 可用性：开通一个账号不超过 3 次操作 \\
\midrule
设计约束 & 复用 UC-08 的登录入口，角色决定登录后的功能范围 \\
\midrule
\bottomrule
\end{longtable}}
\noindent（13）预定商品

{\small
\begin{longtable}{@{}p{2.4cm} p{12.9cm}@{}}
\caption{预定商品用例规约}\label{tab:uc-13}\\
\toprule
\textbf{标题} & \textbf{内容} \\
\midrule
\endfirsthead
\toprule
\textbf{标题} & \textbf{内容} \\
\midrule
\endhead
用例编号：用例名 & UC-13：预定商品 \\
\midrule
执行者 & 客户（无需登录） \\
\midrule
前置条件 & 目标货品处于“在柜”且在展示范围内，且当前没有其他有效预定 \\
\midrule
后置条件 & 新增一条预定记录，关联货品与联系方式；货品状态不变，不锁货 \\
\midrule
涉众利益 & 客户：先把看中的货留住，到店有货可看；\newline 销售：知道哪件货被预定，接待有准备；\newline 老板：预定数据反映真实到店意向，防止恶意占货 \\
\midrule
基本路径 & 1. 客户选择货品并请求预定；\newline 2. 客户提供姓名与手机号；\newline 3. 系统校验该手机号的有效预定次数未达上限，且该货品没有其他有效预定；\newline 4. 系统保存预定记录并生成预定编号；\newline 5. 系统反馈预定结果 \\
\midrule
扩展 & 3a. 该手机号有效预定已达 3 次：3a1 系统反馈需等待已有预定完成后再预定；3a2 返回2；\newline 3b. 该货品已有其他有效预定：3b1 系统反馈货品已被预定；3b2 返回2 \\
\midrule
字段列表 & 预定记录 = 预定编号 + 货品 + 姓名 + 手机号 + 提交时间 + 状态{待到店|已达成|已爽约|已取消} \\
\midrule
业务规则 & ①一个手机号最多同时存在 3 条有效预定；\newline ②预定如约达成（按时到店，由销售确认）后恢复 1 次额度；\newline ③爽约或客户取消不恢复额度；\newline ④同一货品同时只能有 1 条有效预定；\newline ⑤预定超过约定到店时间未到店，系统标记爽约并扣减额度；\newline ⑥客户预定后联系销售并确定到店的，销售将该货品置为“预留”；到店看过暂不决定的，销售取消预留，货品回“在柜” \\
\midrule
质量需求 & 性能：提交后 1 秒内反馈；\newline 可用性：完成一次预定不超过 3 次操作 \\
\midrule
设计约束 & 预定入口从货品展示页进入，剩余预定次数在页面上可见 \\
\midrule
\bottomrule
\end{longtable}}
\noindent（14）提交货品出让申请

{\small
\begin{longtable}{@{}p{2.4cm} p{12.9cm}@{}}
\caption{提交货品出让申请用例规约}\label{tab:uc-14}\\
\toprule
\textbf{标题} & \textbf{内容} \\
\midrule
\endfirsthead
\toprule
\textbf{标题} & \textbf{内容} \\
\midrule
\endhead
用例编号：用例名 & UC-14：提交货品出让申请 \\
\midrule
执行者 & 供应商（含玉料供应商、成品供应商、回流出让方，账号由老板开通） \\
\midrule
前置条件 & 无（身份验证由 UC-08 提供） \\
\midrule
后置条件 & 新增一条出让申请，状态“待审核”，等待老板审核 \\
\midrule
涉众利益 & 供应商：在线提交货品信息，不用上门先谈一轮；\newline 老板：申请集中在线上，有空再审核，货源线索不散失；\newline 销售：审核通过的货源到店后按常规流程接待 \\
\midrule
基本路径 & 1. 供应商选择出让类型（玉料/成品/回流）并填写货品信息（品名、数量、期望出让价、描述），提供料子图片；\newline 2. 系统校验必填字段完整；\newline 3. 系统保存申请并生成申请编号，状态置为“待审核”；\newline 4. 系统反馈提交结果 \\
\midrule
扩展 & 2a. 必填字段缺失：2a1 系统反馈应补全后提交；2a2 返回1 \\
\midrule
字段列表 & 出让申请 = 申请编号 + 供应商 + 出让类型{玉料|成品|回流} + 品名 + 数量 + 期望出让价 + 描述 + 料子图片 + 提交时间 + 状态{待审核|已通过|已驳回} \\
\midrule
业务规则 & ①玉料供应商、成品供应商与回流出让方共用同一申请入口，类型取自供应商档案；\newline ②申请经老板审核通过后，货品实际到店时按 UC-02 建档入库；\newline ③期望出让价仅供审核参考，不对外展示 \\
\midrule
质量需求 & 性能：提交后 1 秒内反馈；\newline 可用性：完成一次出让申请不超过 5 次操作 \\
\midrule
设计约束 & 申请入口在供应商登录后的界面，与调货价目查询互不可见 \\
\midrule
\bottomrule
\end{longtable}}
\noindent（15）预约上门时间

{\small
\begin{longtable}{@{}p{2.4cm} p{12.9cm}@{}}
\caption{预约上门时间用例规约}\label{tab:uc-15}\\
\toprule
\textbf{标题} & \textbf{内容} \\
\midrule
\endfirsthead
\toprule
\textbf{标题} & \textbf{内容} \\
\midrule
\endhead
用例编号：用例名 & UC-15：预约上门时间 \\
\midrule
执行者 & 供应商、加工师傅（账号由老板开通） \\
\midrule
前置条件 & 供应商有已通过的出让申请，或加工师傅有已确认的完工交付申请 \\
\midrule
后置条件 & 新增一条上门预约，关联预约方、事项与时段 \\
\midrule
涉众利益 & 供应商与加工师傅：约好时间再来，不白跑；\newline 老板：交接时间集中安排，减少临时接待；\newline 销售：知道什么时间有什么交接，提前准备 \\
\midrule
基本路径 & 1. 预约方选择要办理的事项（出让交货/完工交付）；\newline 2. 预约方提供期望上门时段；\newline 3. 系统校验该时段没有其他上门预约冲突，保存预约；\newline 4. 系统反馈预约结果 \\
\midrule
扩展 & 3a. 时段已有其他上门预约：3a1 系统反馈应更换时段；3a2 返回2 \\
\midrule
字段列表 & 上门预约 = 预约编号 + 预约方 + 事项 + 期望时段 + 提交时间 + 状态{待上门|已完成|已取消} \\
\midrule
业务规则 & ①预约事项必须关联一条已通过的出让申请或已确认的完工交付申请；\newline ②同一时段只安排一次上门交接；\newline ③预约由老板在系统内确认，改期需重新预约 \\
\midrule
质量需求 & 性能：提交后 1 秒内反馈；\newline 可用性：预约不超过 3 次操作 \\
\midrule
设计约束 & 上门交接与客户到店预约统一安排时段，避免相互冲突 \\
\midrule
\bottomrule
\end{longtable}}
\noindent（16）查看加工任务

{\small
\begin{longtable}{@{}p{2.4cm} p{12.9cm}@{}}
\caption{查看加工任务用例规约}\label{tab:uc-16}\\
\toprule
\textbf{标题} & \textbf{内容} \\
\midrule
\endfirsthead
\toprule
\textbf{标题} & \textbf{内容} \\
\midrule
\endhead
用例编号：用例名 & UC-16：查看加工任务 \\
\midrule
执行者 & 加工师傅（账号由老板开通） \\
\midrule
前置条件 & 无（身份验证由 UC-08 提供） \\
\midrule
后置条件 & 只读查询，不改变任何数据 \\
\midrule
涉众利益 & 加工师傅：随时掌握待加工与在加工任务、要求与工费；\newline 老板：不用电话逐个通知师傅；\newline 销售：加工进度可查，客户问货时能答复 \\
\midrule
基本路径 & 1. 加工师傅请求查看加工任务；\newline 2. 系统检索分派给该师傅的任务（货品、加工要求、工费、期限、状态）；\newline 3. 系统反馈查询结果 \\
\midrule
扩展 & 2a. 没有任何任务：2a1 系统反馈当前无加工任务 \\
\midrule
字段列表 & 加工任务 = 任务编号 + 货品 + 加工要求 + 工费 + 期限 + 状态{待加工|在加工|待交付|已交付} \\
\midrule
业务规则 & ①任务由老板在安排加工时生成，加工师傅只能查看，不能改任务内容；\newline ②工费金额对师傅可见，结算仍由老板在系统内记录 \\
\midrule
质量需求 & 性能：查询 2 秒内反馈；\newline 可靠性：任务状态与老板端记录一致 \\
\midrule
设计约束 & 加工师傅与内部账号共用登录入口，登录后只见任务相关功能 \\
\midrule
\bottomrule
\end{longtable}}
\noindent（17）提交完工交付申请

{\small
\begin{longtable}{@{}p{2.4cm} p{12.9cm}@{}}
\caption{提交完工交付申请用例规约}\label{tab:uc-17}\\
\toprule
\textbf{标题} & \textbf{内容} \\
\midrule
\endfirsthead
\toprule
\textbf{标题} & \textbf{内容} \\
\midrule
\endhead
用例编号：用例名 & UC-17：提交完工交付申请 \\
\midrule
执行者 & 加工师傅（账号由老板开通） \\
\midrule
前置条件 & 存在状态为“在加工”的本人任务 \\
\midrule
后置条件 & 任务状态改为“待交付”，等待老板确认并安排入库 \\
\midrule
涉众利益 & 加工师傅：完工即在线报备，不用等老板有空再交接；\newline 老板：完工信息第一时间可见，及时安排交接入库；\newline 销售：货品加工完成的信息可查，便于对客户答复 \\
\midrule
基本路径 & 1. 加工师傅选择已完工的加工任务；\newline 2. 加工师傅提供加工前、加工中与完工的实物图与完工说明；\newline 3. 系统校验任务状态为“在加工”，更新为“待交付”；\newline 4. 系统反馈提交结果 \\
\midrule
扩展 & 3a. 任务状态不是“在加工”：3a1 系统反馈不允许重复提交；3a2 返回1 \\
\midrule
字段列表 & 完工交付申请 = 申请编号 + 任务编号 + 完工说明 + 过程图（加工前/加工中/完工）+ 提交时间 + 状态{待确认|已确认} \\
\midrule
业务规则 & ①完工交付申请提交后，货品入库仍由老板按 UC-02 建档，加工师傅不录入货品档案；\newline ②老板确认完工后，师傅通过 UC-15 预约上门时间交货；\newline ③一条任务只能有一条待确认的完工交付申请 \\
\midrule
质量需求 & 性能：提交后 1 秒内反馈；\newline 可靠性：完工状态与老板端一致 \\
\midrule
设计约束 & 提交入口从加工任务列表进入；老板确认后按 UC-15 预约交接时段 \\
\midrule
\bottomrule
\end{longtable}}
\noindent（18）审核出让申请

{\small
\begin{longtable}{@{}p{2.4cm} p{12.9cm}@{}}
\caption{审核出让申请用例规约}\label{tab:uc-18}\\
\toprule
\textbf{标题} & \textbf{内容} \\
\midrule
\endfirsthead
\toprule
\textbf{标题} & \textbf{内容} \\
\midrule
\endhead
用例编号：用例名 & UC-18：审核出让申请 \\
\midrule
执行者 & 老板（主执行者） \\
\midrule
前置条件 & 存在状态为“待审核”的出让申请 \\
\midrule
后置条件 & 申请状态改为“已通过”或“已驳回”；通过的申请在货品到店时转入入库流程 \\
\midrule
涉众利益 & 老板：货源线索集中审核，通过与否自己把关；\newline 供应商：审核结果在线可查，不用反复电话确认；\newline 销售：通过的货源到店后按常规流程接待 \\
\midrule
基本路径 & 1. 老板请求查看待审核的出让申请；\newline 2. 系统反馈申请详情与期望出让价；\newline 3. 老板选择通过或驳回（驳回须填写原因）；\newline 4. 系统保存审核结果并反馈 \\
\midrule
扩展 & 3a. 驳回未填原因：3a1 系统反馈应填写原因后提交；3a2 返回3 \\
\midrule
字段列表 & 出让申请（含审核） = 出让申请字段 + 审核意见 + 审核时间 + 审核人 \\
\midrule
业务规则 & ①审核通过的申请，货品实际到店时按 UC-02 建档，来源类型与出让类型一致；\newline ②驳回的申请保留记录，不删除；\newline ③审核只能由老板执行，销售可查看但不能审核 \\
\midrule
质量需求 & 性能：审核提交后 1 秒内反馈 \\
\midrule
设计约束 & 审核入口在老板登录后的界面，待审核列表按提交时间排序；通过后由供应商按 UC-15 预约上门交接 \\
\midrule
\bottomrule
\end{longtable}}

\noindent（19）登记货款记录

{\small
\begin{longtable}{@{}p{2.4cm} p{12.9cm}@{}}
\caption{登记货款记录用例规约}\label{tab:uc-19}\\
\toprule
\textbf{标题} & \textbf{内容} \\
\midrule
\endfirsthead
\toprule
\textbf{标题} & \textbf{内容} \\
\midrule
\endhead
用例编号：用例名 & UC-19：登记货款记录 \\
\midrule
执行者 & 老板（主执行者） \\
\midrule
前置条件 & 存在对应业务单据（采购入库、调货往来或加工任务） \\
\midrule
后置条件 & 新增一条货款记录（业务类型、关联单据、金额、收付方向、经手人），货款往来可追溯 \\
\midrule
涉众利益 & 老板：每笔货款可追溯到来源单据与经手人，资金往来一本账；\newline
供应商与加工师傅：货款收付有据可查，结算不扯皮；\newline
销售：不承担记账职责 \\
\midrule
基本路径 & 1. 老板选择业务类型（采购货款/调货货款/加工工费/寄售货款）并关联对应单据；\newline
2. 老板录入金额与收付方向（支付/收取）；\newline
3. 系统校验金额与单据一致；\newline
4. 系统保存货款记录；\newline
5. 系统反馈登记结果 \\
\midrule
扩展 & 3a. 金额与单据不一致：3a1 系统反馈应核对后重新录入；3a2 返回2 \\
\midrule
字段列表 & 货款记录 = 记录编号 + 业务类型{采购货款|调货货款|加工工费|寄售货款} + 关联单据 + 金额 + 收付方向{支付|收取} + 经手人 + 时间 \\
\midrule
业务规则 & ①每条货款记录必须关联来源单据；\newline ②收付方向区分支付与收取；\newline ③记录只增不改，登记错误时作废重录 \\
\midrule
质量需求 & 可用性：登记一条不超过 3 次操作 \\
\midrule
设计约束 & 与统计口径衔接，付款计入成本侧、收款计入销售侧 \\
\midrule
\bottomrule
\end{longtable}}

\noindent（20）维护调货记录

{\small
\begin{longtable}{@{}p{2.4cm} p{12.9cm}@{}}
\caption{维护调货记录用例规约}\label{tab:uc-20}\\
\toprule
\textbf{标题} & \textbf{内容} \\
\midrule
\endfirsthead
\toprule
\textbf{标题} & \textbf{内容} \\
\midrule
\endhead
用例编号：用例名 & UC-20：维护调货记录 \\
\midrule
执行者 & 老板（主执行者） \\
\midrule
前置条件 & 存在调货业务（合作商户调入或调出） \\
\midrule
后置条件 & 新增或更新调货往来记录与调货库存记录，调货货品去向可追溯 \\
\midrule
涉众利益 & 老板：调货往来与库存一本账，结算不扯皮；\newline
合作商户：往来记录准确，对账有依据；\newline
销售：库存状态及时更新，接待有准备 \\
\midrule
基本路径 & 1. 老板选择调货业务（调入/调出）并确认货品与数量；\newline
2. 系统新增调货往来记录；\newline
3. 老板登记货品去向（售出/退回），系统更新调货库存记录；\newline
4. 系统反馈登记结果 \\
\midrule
扩展 & 3a. 货品状态与调货记录不符：3a1 系统反馈应核对后重新登记；3a2 返回3 \\
\midrule
字段列表 & 调货记录 = 记录编号 + 方向{调入|调出} + 合作商户 + 货品 + 数量 + 状态{在调|已售|已退回} + 时间 \\
\midrule
业务规则 & ①往来记录与库存记录联动更新；\newline ②调给同行的货以客户鉴赏期（3 天或 7 天）届满为结账时点；\newline ③鉴赏期内客户确定留货即按已售出处理，触发与商户的货款结算 \\
\midrule
质量需求 & 可靠性：往来记录与老板端一致 \\
\midrule
设计约束 & 调货记录仅老板可维护，销售只读 \\
\midrule
\bottomrule
\end{longtable}}

\noindent（21）登记寄售合同

{\small
\begin{longtable}{@{}p{2.4cm} p{12.9cm}@{}}
\caption{登记寄售合同用例规约}\label{tab:uc-21}\\
\toprule
\textbf{标题} & \textbf{内容} \\
\midrule
\endfirsthead
\toprule
\textbf{标题} & \textbf{内容} \\
\midrule
\endhead
用例编号：用例名 & UC-21：登记寄售合同 \\
\midrule
执行者 & 老板（主执行者） \\
\midrule
前置条件 & 供应商在基础数据中已建档 \\
\midrule
后置条件 & 新增一份状态为“生效”的寄售合同与寄售明细，货品以寄售方式入库（来源类型为寄售、成本为寄售价、状态为在柜） \\
\midrule
涉众利益 & 老板：寄售货品有合同与寄售价为据，售出后结算不扯皮；\newline
供应商：货品在店期间去向可查，售出按寄售价回款；\newline
销售：寄售货品与自有货品同场展示，无须区分维护 \\
\midrule
基本路径 & 1. 老板选择供应商并填写合同信息（到期日、结算方式）；\newline
2. 系统校验供应商存在，创建寄售合同；\newline
3. 老板逐件录入寄售货品（品名、类别、克重与寄售价）；\newline
4. 系统按寄售方式建档入库并创建寄售明细；\newline
5. 系统写入库存流水；\newline
6. 系统反馈登记结果（含合同编号） \\
\midrule
扩展 & 2a. 供应商不存在：2a1 系统反馈先在维护基础数据中建档；2a2 返回1 \\
\midrule
字段列表 & 寄售合同 = 合同编号 + 供应商 + 签订时间 + 到期日 + 结算方式{售出即结|按月结算} + 状态{生效|已终止|已到期}；\newline 寄售明细 = 合同 + 货品 + 寄售价 + 结算状态{未结算|已结算} \\
\midrule
业务规则 & ①一份合同可含多件寄售货品，一件货品同一时间只挂一份未结清合同；\newline ②寄售货品的成本记为寄售价，作为结算依据；\newline ③合同终止或到期后，未结算明细对应的货品退回供应商 \\
\midrule
质量需求 & 可用性：登记一件寄售货品不超过 4 次操作 \\
\midrule
设计约束 & 寄售货品与自有货品在柜面同场展示，售出流程与 UC-01 相同 \\
\midrule
\bottomrule
\end{longtable}}

\noindent（22）结算寄售货款

{\small
\begin{longtable}{@{}p{2.4cm} p{12.9cm}@{}}
\caption{结算寄售货款用例规约}\label{tab:uc-22}\\
\toprule
\textbf{标题} & \textbf{内容} \\
\midrule
\endfirsthead
\toprule
\textbf{标题} & \textbf{内容} \\
\midrule
\endhead
用例编号：用例名 & UC-22：结算寄售货款 \\
\midrule
执行者 & 老板（主执行者） \\
\midrule
前置条件 & 存在已售出且未结算的寄售明细 \\
\midrule
后置条件 & 新增一条台账结算（类型为寄售），对应寄售明细置为已结算 \\
\midrule
涉众利益 & 老板：按寄售价与供应商结清，账实一致；\newline
供应商：回款有结算单为据；\newline
销售：不承担结算职责 \\
\midrule
基本路径 & 1. 老板选择已售出的寄售货品；\newline
2. 系统校验该明细未结算；\newline
3. 系统按寄售价创建台账结算（往来方为供应商）；\newline
4. 系统将寄售明细置为已结算；\newline
5. 系统反馈结算结果（含结算编号） \\
\midrule
扩展 & 2a. 明细已结算：2a1 系统反馈该货品已结清；2a2 返回1 \\
\midrule
字段列表 & 台账结算 = 结算编号 + 往来方 + 结算类型{调货|寄售} + 结算金额 + 结算方式{现金|转账|冲抵} + 结算时间 + 经手人 \\
\midrule
业务规则 & ①结算金额为寄售价；\newline ②一张结算单只对应一张业务单据（寄售明细或调货记录）；\newline ③结算完成后资金流水在 UC-19 货款记录登记 \\
\midrule
质量需求 & 可用性：结算一笔不超过 3 次操作 \\
\midrule
设计约束 & 寄售与调货结算共用一本台账 \\
\midrule
\bottomrule
\end{longtable}}

\noindent（23）结算调货台账

{\small
\begin{longtable}{@{}p{2.4cm} p{12.9cm}@{}}
\caption{结算调货台账用例规约}\label{tab:uc-23}\\
\toprule
\textbf{标题} & \textbf{内容} \\
\midrule
\endfirsthead
\toprule
\textbf{标题} & \textbf{内容} \\
\midrule
\endhead
用例编号：用例名 & UC-23：结算调货台账 \\
\midrule
执行者 & 老板（主执行者） \\
\midrule
前置条件 & 存在未结算的调货记录 \\
\midrule
后置条件 & 新增一条台账结算（类型为调货），所选调货记录置为已结算 \\
\midrule
涉众利益 & 老板：按商户对账，鉴赏期届满即结清；\newline
合作商户：往来与结算有据可查，不扯皮；\newline
销售：不承担对账职责 \\
\midrule
基本路径 & 1. 老板选择合作商户，系统列出未结算调货记录；\newline
2. 老板勾选本次结算的记录并确认结算方式；\newline
3. 系统汇总金额并创建台账结算；\newline
4. 系统将所选调货记录置为已结算；\newline
5. 系统反馈结算结果（含结算编号） \\
\midrule
扩展 & 1a. 无未结算记录：1a1 系统反馈该商户往来已结清；1a2 用例结束 \\
\midrule
字段列表 & 同台账结算（结算类型为调货），金额为所选调货记录的汇总 \\
\midrule
业务规则 & ①结算金额为所选调货记录金额汇总；\newline ②已结算的调货记录不再出现在结算清单；\newline ③与 UC-20 规则②③衔接，鉴赏期届满为结账时点 \\
\midrule
质量需求 & 可靠性：结算后单据与商户端记录一致 \\
\midrule
设计约束 & 结算清单支持一次多笔 \\
\midrule
\bottomrule
\end{longtable}}

\noindent（24）登记分期计划

{\small
\begin{longtable}{@{}p{2.4cm} p{12.9cm}@{}}
\caption{登记分期计划用例规约}\label{tab:uc-24}\\
\toprule
\textbf{标题} & \textbf{内容} \\
\midrule
\endfirsthead
\toprule
\textbf{标题} & \textbf{内容} \\
\midrule
\endhead
用例编号：用例名 & UC-24：登记分期计划 \\
\midrule
执行者 & 销售（主执行者） \\
\midrule
前置条件 & 存在成交记录，且该成交未关联分期计划 \\
\midrule
后置条件 & 新增一份状态为“未还”的分期计划，计划总额等于成交价 \\
\midrule
涉众利益 & 销售：大额成交当场可谈分期，不丢单；\newline
客户：分期减轻一次付清压力；\newline
老板：分期账目清楚，逾期可追踪 \\
\midrule
基本路径 & 1. 销售选择成交记录并填写分期期数；\newline
2. 系统校验该成交未关联分期计划；\newline
3. 系统按成交价创建分期计划（每期应还、下一期到期日、状态未还）；\newline
4. 系统反馈登记结果（含计划编号） \\
\midrule
扩展 & 2a. 已存在分期计划：2a1 系统反馈一份成交只有一份计划；2a2 用例结束 \\
\midrule
字段列表 & 分期计划 = 计划编号 + 成交记录 + 分期期数 + 每期应还 + 已还金额 + 下一期到期日 + 状态{未还|逾期|已结清} \\
\midrule
业务规则 & ①一份成交至多一份分期计划，计划总额等于成交价；\newline ②分期期间货款按已收金额累计，不影响货品所有权转移 \\
\midrule
质量需求 & 可用性：登记一份计划不超过 3 次操作 \\
\midrule
设计约束 & 分期期数与门槛由店内规则决定，系统不设上限 \\
\midrule
\bottomrule
\end{longtable}}

\noindent（25）记录分期还款

{\small
\begin{longtable}{@{}p{2.4cm} p{12.9cm}@{}}
\caption{记录分期还款用例规约}\label{tab:uc-25}\\
\toprule
\textbf{标题} & \textbf{内容} \\
\midrule
\endfirsthead
\toprule
\textbf{标题} & \textbf{内容} \\
\midrule
\endhead
用例编号：用例名 & UC-25：记录分期还款 \\
\midrule
执行者 & 销售（主执行者） \\
\midrule
前置条件 & 存在未结清的分期计划 \\
\midrule
后置条件 & 更新分期计划的已还金额与下一期到期日，结清后状态置为已结清 \\
\midrule
涉众利益 & 客户：还款进度随时可查；\newline
销售：收款当场登记，不用事后补账；\newline
老板：逾期计划一目了然 \\
\midrule
基本路径 & 1. 销售选择分期计划并录入还款金额；\newline
2. 系统校验金额不超过未还余额；\newline
3. 系统累加已还金额并推进下一期到期日；\newline
4. 系统校验是否结清，结清则置为已结清；\newline
5. 系统反馈还款结果（含剩余应还） \\
\midrule
扩展 & 2a. 金额超过未还余额：2a1 系统反馈应按余额录入；2a2 返回1 \\
\midrule
字段列表 & 同分期计划 \\
\midrule
业务规则 & ①到期日过仍有未还即置为逾期；\newline ②还款可跨期一次结清；\newline ③结清后计划归档，不再接受还款 \\
\midrule
质量需求 & 可用性：记录一次还款不超过 2 次操作 \\
\midrule
设计约束 & 逾期状态在查询时按到期日与未还余额实时判定（超过到期日未还即逾期） \\
\midrule
\bottomrule
\end{longtable}}

\noindent（26）查看库存流水

{\small
\begin{longtable}{@{}p{2.4cm} p{12.9cm}@{}}
\caption{查看库存流用例规约}\label{tab:uc-26}\\
\toprule
\textbf{标题} & \textbf{内容} \\
\midrule
\endfirsthead
\toprule
\textbf{标题} & \textbf{内容} \\
\midrule
\endhead
用例编号：用例名 & UC-26：查看库存流水 \\
\midrule
执行者 & 老板（主执行者）；销售、合作商户（只读） \\
\midrule
前置条件 & 无（登录即可） \\
\midrule
后置条件 & 无数据变更，返回满足条件的库存流水清单 \\
\midrule
涉众利益 & 老板：任一货品的进出有账可查，调货库存一本账；\newline
销售：交接班时核对货品状态；\newline
合作商户：对账时可当场出示流水 \\
\midrule
基本路径 & 1. 老板选择货品或时间段；\newline
2. 系统按条件查询库存流水；\newline
3. 系统返回流水清单（变更前状态、变更后状态、关联单据、经手人、时间） \\
\midrule
扩展 & 2a. 无满足条件的流水：2a1 系统反馈该条件下无记录 \\
\midrule
字段列表 & 库存流水 = 货品 + 变更前状态 + 变更后状态 + 关联单据 + 经手人 + 流水时间 \\
\midrule
业务规则 & ①货品每次状态变更必须写一条流水；\newline ②流水写入后不可修改、不可删除 \\
\midrule
质量需求 & 性能：单货品流水查询 3 秒内返回 \\
\midrule
设计约束 & 流水状态取值与货品状态机（07a）一致 \\
\midrule
\bottomrule
\end{longtable}}

% ==================== C 分析（承接上文，无过渡） ====================

\newpage

% ============================================================
% C 阶段：分析
% ============================================================
\section{C 分析}

\subsection{分析类图（边界 / 控制 / 实体三分法）}

分析类图把系统涉及的类按职责分成三类：
\textbf{边界类}是系统与外部系统交互的界面，一个有接口的外部系统映射一个边界类；
只负责收发消息、不做业务判断；\textbf{控制类}负责协调与流程控制，把执行者请求分配给实体类完成，
不持有业务数据，因此不写 get/set 一类的存取操作；\textbf{实体类}是核心域概念，
属性只写直接描述该类特征的领域字段，不画标识属性（主键）与状态属性（状态交由状态机负责）。

\subsubsection{边界类与控制类}

边界类与控制类分两张图给出：柜台电脑端一侧见图~\ref{fig:05a-boundary-control-class}，微信小程序一侧见图~\ref{fig:05a2-boundary-control-class}——一个外部系统一张图。边界类按\textbf{外部系统}划分而不是按角色划分：
柜台电脑端（老板与销售共用一套系统）与微信小程序（客户、供应商、加工师傅、合作商户共用一个小程序）
各对应一个边界类，合计 2 个；同一外部系统内各角色能看到的字段与按钮由登录角色决定，不另立边界类。
登录在两个外部系统都有入口，因此登录控制被两个边界类共同依赖。
控制类与阶段 B 的用例一一对应，命名一律“用例名 + 控制”，每个控制类只承担“接收请求、组织协作、返回结果”的责任。
边界类到控制类是转发关系，控制类到实体类是分配责任，前者见表~\ref{tab:b2c}，后者见各用例的分析序列图。
两者的方向都是\textbf{依赖}：控制类不保存数据，也不调用实体类的取值/赋值操作。

\fig{分析类图（一）：边界类与柜台端控制类}{05a-boundary-control-class}

\fig{分析类图（二）：微信小程序边界类与控制类}{05a2-boundary-control-class}

\needspace{8\baselineskip}
\begin{table}[H]
\centering
\small
\caption{边界类到控制类的对应关系（依赖）}
\label{tab:b2c}
\begin{tabularx}{\textwidth}{@{}l X@{}}
\toprule
边界类 & 覆盖的控制类（对应用例） \\
\midrule
柜台端界面 & 登记成交、登记货品入库、登记退货、维护客户信息、登记送检与证书、查看销售与周转统计、维护基础数据、开通外部账户、审核出让申请、登记货款记录、维护调货记录、登记寄售合同、结算寄售货款、结算调货台账、登记分期计划、记录分期还款、查看库存流水（共 17 个） \\
微信小程序界面 & 浏览货品、预约查看货品、查看调货价格、预定商品、提交货品出让申请、预约上门时间、查看加工任务、提交完工交付申请（共 8 个） \\
两端共用 & 登录（UC-08 被各端用例以 include 引用，在两个外部系统都有入口） \\
\bottomrule
\end{tabularx}
\end{table}


\subsubsection{实体类及其关联（关联与组合）}

实体类及其关联如图~\ref{fig:05b1-entity-goods}。共 26 个实体类：货品、入库批次、证书、成交、退货、
送检、预定、预约、出让申请、上门预约、加工任务、完工交付申请、货款、调货、寄售合同、寄售明细、
台账结算、分期计划、库存流水、客户、来源方、鉴定机构、货品类别、用户、角色、权限。
来源方是供应商、加工师傅与同行商户的合一：三类的属性与操作完全一致，只以类别取值区分
（与货款按业务类型取值区分同理），类别在同一取值内名称唯一。
实体类的操作列与下节分析序列图的消息一一对应——某条消息传进来，就是该实体类的操作。
类名不带“记录/信息/数据/情况”等后缀，不画标识属性（主键），
也不画名为“状态”的属性（货品状态交由状态机负责）。合同编号、证书编号等业务编码是印在纸质单据上、用户可引用的领域字段，不是主键标识，故保留在类图中；主键（代理键）不画，数据层与业务编码分离。

客户类有一条自反关联：来源渠道含“朋友介绍”，因而客户之间存在介绍关系。
\textbf{普通关联}用于可共享、或同一对对象可多次发生的关系（如同一件货品的多次成交）；\textbf{组合}用于部分不可共享、且不能脱离整体存在的关系，共四处——证书挂货品、完工交付申请挂加工任务、货品入库批次、寄售明细挂寄售合同；控制类到实体类的\textbf{依赖}见分析类图（一）。\textbf{泛化}不采用：各类账户的差别只在权限与登录入口，核心属性完全一致，属于属性值差别；用户与角色为多对多、角色与权限也为多对多，角色与权限各自成类成表，不为五类角色各建子类。
类之间的 33 条关联连同多重性与角色名见表~\ref{tab:rels}。

\fig{实体类之间的关联（一）：货品与入库批次、证书}{05b1-entity-goods}
\fig{实体类之间的关联（二）：货品与成交}{05b2-entity-deal}
\figtwoW{0.72}{实体类之间的关联（三）：送检与证书}{05b3-entity-inspect}{1.0}{实体类之间的关联（四）：预约、货品、预定与客户}{05b4-entity-book}
\figtwo[0.60]{实体类之间的关联（五）：来源方、出让申请与上门预约}{05b5-entity-transfer}{实体类之间的关联（六）：货品、加工任务与来源方}{05b6-entity-process}
\fig{实体类之间的关联（七）：货品类别与货品的分类关系}{05b8-entity-category}
\fig{实体类之间的关联（八）：货品送检与送检单}{05b9-entity-goods-inspect}
\fig{实体类之间的关联（九）：调货与来源方往来}{05b11-entity-money-partner}
\figtwo[0.72]{实体类之间的关联（十）：加工任务与完工交付申请}{05b12-entity-delivery}{实体类之间的关联（十一）：送检、鉴定机构与证书}{05b13-entity-cert-org}
\fig{实体类之间的关联（十二）：货品的调拨}{05b14-entity-goods-transfer2}
\figtwo[0.72]{实体类之间的关联（十三）：成交与客户下单}{05b15-entity-deal-customer}{实体类之间的关联（十四）：成交与退货}{05b16-entity-deal-return}
\fig{实体类之间的关联（十五）：来源方、寄售合同、寄售明细与货品}{05b17-entity-consign}
\fig{实体类之间的关联（十六）：来源方、台账结算、调货与寄售明细}{05b18-entity-settle}
\fig{实体类之间的关联（十七）：成交与分期计划}{05b19-entity-installment}
\fig{实体类之间的关联（十八）：货品、用户与库存流水}{05b20-entity-stocklog}
\fig{实体类之间的关联（十九）：用户、角色与权限}{05b21-entity-rbac}

\needspace{8\baselineskip}
\begin{table}[H]
\centering
\small
\caption{实体类之间的关联（多重性与角色名）}
\label{tab:rels}
\begin{tabularx}{\textwidth}{@{}l c l c X@{}}
\toprule
源类 & 源端 & 目标类 & 目标端 & 角色名与说明 \\
\midrule
入库批次 & 1 & 货品 & 0..* & 含货品（一件货品只属一个批次）\textbf{（组合）} \\
货品 & 1 & 证书 & 0..* & 挂证书（证书不可共享，货品注销则证书无意义）\textbf{（组合）} \\
加工任务 & 1 & 完工交付申请 & 0..* & 产出（完工交付申请不可脱离任务）\textbf{（组合）} \\
来源方 & 1 & 寄售合同 & 0..* & 寄出（类别=供应商） \\
寄售合同 & 1 & 寄售明细 & 0..* & 含（一份合同多件寄售货品）\textbf{（组合）} \\
货品 & 1 & 寄售明细 & 0..* & 被寄售（一货品同一时间至多一条未结清明细，结算后可再寄售） \\
来源方 & 1 & 台账结算 & 0..* & 结算往来 \\
调货 & 1 & 台账结算 & 0..1 & 结算目标（一笔调货至多一次结算） \\
寄售明细 & 1 & 台账结算 & 0..1 & 结算目标（一件寄售货品至多一次结算） \\
成交 & 1 & 分期计划 & 0..1 & 分期（一份成交至多一份计划） \\
货品 & 1 & 库存流水 & 0..* & 记录（每次状态变更必写流水） \\
用户 & 1 & 库存流水 & 0..* & 经手 \\
用户 & 0..* & 角色 & 0..* & 具有（多对多，落表 t\_user\_role） \\
角色 & 0..* & 权限 & 0..* & 授予（多对多，落表 t\_role\_permission） \\
货品类别 & 1 & 货品 & 0..* & 分类 \\
货品 & 1 & 成交 & 0..* & 售出（退货后仍可再售，故可多次） \\
客户 & 0..1 & 成交 & 0..* & 下单（散客不关联客户档案，允许为空） \\
成交 & 1 & 退货 & 0..* & 被退（退货是后续事件而非组成部分，故不画组合） \\
货品 & 1 & 送检 & 0..* & 送检 \\
送检 & 0..1 & 证书 & 0..* & 出具（证书由送检出具，随货品带入的证书无对应送检） \\
鉴定机构 & 1 & 送检 & 0..* & 送检机构 \\
鉴定机构 & 1 & 证书 & 0..* & 签发机构 \\
货品 & 1 & 预定 & 0..* & 被预定 \\
预定 & 0..* & 客户 & 0..1 & 登记给（预定以姓名与手机号提交，已建档客户才关联） \\
客户 & 0..1 & 客户 & 0..* & 介绍（自反关联：来源渠道含“朋友介绍”） \\
货品 & 1 & 预约 & 0..* & 被预约看 \\
来源方 & 1 & 出让申请 & 0..* & 提交（类别=供应商，含回流出让方） \\
出让申请 & 0..1 & 上门预约 & 0..* & 预约上门 \\
完工交付申请 & 0..1 & 上门预约 & 0..* & 预约上门（UC-15 规则①） \\
货品 & 1 & 加工任务 & 0..* & 加工 \\
来源方 & 1 & 加工任务 & 0..* & 承接（类别=加工师傅） \\
货品 & 1 & 调货 & 0..* & 调拨 \\
来源方 & 1 & 调货 & 0..* & 往来方（类别=同行商户） \\
\bottomrule
\end{tabularx}
\end{table}

\subsection{分析序列图}

序列图对应每个用例的路径，每个用例的基本路径各画一张图：
简单的扩展点合并到基本路径图，复杂扩展点单独一张并与基本路径图链接。
研究对象是系统内各类，软件类——执行者只发消息到边界类，
边界类把请求转给控制类，控制类再组织实体类完成业务。
消息代表责任分配而非数据流动，控制类只分配责任，不取存实体类的属性。
26 个用例的基本路径逐一给出，如图~\ref{fig:06a-uc01-sale} 至图~\ref{fig:06y-uc26-stock-log}。

\fig[1.0]{登记成交（UC-01）分析序列图}{06a-uc01-sale}
\fig[1.0]{登记货品入库（UC-02）分析序列图}{06b-uc02-goods-in}
\fig[1.0]{登记退货（UC-03）分析序列图}{06c-uc03-return}
\fig[1.0]{维护客户信息（UC-04）分析序列图}{06i-uc04-customer}
\fig[1.0]{登记送检与证书（UC-05）分析序列图}{06d-uc05-inspect}
\fig[1.0]{查看销售与周转统计（UC-06）分析序列图}{06j-uc06-stats}
\fig[1.0]{维护基础数据（UC-07）分析序列图}{06k-uc07-basicdata}
\fig[1.0]{登录（UC-08）分析序列图}{06l-uc08-login}
\fig[1.0]{浏览货品（UC-09）分析序列图}{06m-uc09-browse}
\fig[1.0]{预约查看货品（UC-10）分析序列图}{06n-uc10-reserve}
\fig[1.0]{查看调货价格（UC-11）分析序列图}{06o-uc11-transfer-price}
\fig[1.0]{开通外部账户（UC-12）分析序列图}{06p-uc12-open-account}
\fig[1.0]{预定商品（UC-13）分析序列图}{06e-uc13-book}
\fig[1.0]{提交货品出让申请（UC-14）分析序列图}{06f-uc14-transfer-apply}
\fig[1.0]{预约上门时间（UC-15）分析序列图}{06q-uc15-visit-appt}
\fig[1.0]{查看加工任务（UC-16）分析序列图}{06r-uc16-process-task}
\fig[1.0]{提交完工交付申请（UC-17）分析序列图}{06s-uc17-delivery}
\fig[1.0]{审核出让申请（UC-18）分析序列图}{06g-uc18-transfer-review}
\fig[1.0]{登记货款记录（UC-19）分析序列图}{06t-uc19-payment}
\fig[1.0]{维护调货记录（UC-20）分析序列图}{06h-uc20-transfer-rec}
\fig[1.0]{登记寄售合同（UC-21）分析序列图}{06t-uc21-consign}
\fig[1.0]{结算寄售货款（UC-22）分析序列图}{06u-uc22-consign-settle}
\fig[1.0]{结算调货台账（UC-23）分析序列图}{06v-uc23-transfer-settle}
\fig[1.0]{登记分期计划（UC-24）分析序列图}{06w-uc24-installment}
\fig[1.0]{记录分期还款（UC-25）分析序列图}{06x-uc25-installment-pay}
\fig[1.0]{查看库存流水（UC-26）分析序列图}{06y-uc26-stock-log}

\subsection{分析状态机图（只画核心类）}

状态机图只画核心类，每条迁移的触发事件与规约来源一并标注。
共画九张：货品的八状态总览及分主题三张（售出与预留、调货与代销、送检与处置）、出让申请（三状态）、加工任务（四状态）、预定记录（四状态）、寄售合同（三状态）、分期计划（三状态）。
选取依据是“核心类、状态多的事物”：货品是全系统流转最频繁的对象，加工任务是组合结构的根，
预定记录的状态直接驱动额度规则；寄售合同与分期计划是新扩展业务的生命周期载体。

\subsubsection{货品生命周期（八状态，分三张给出）}

货品是全系统流转最频繁的对象，八状态取自 UC-01 的状态枚举：在柜、已订、已售、预留、已调出、代销中、已送检、待处理。
状态之间的迁移共 18 条，总体结构见图~\ref{fig:07a-goods-state}；
再按业务主题分三张细化：售出与预留（图~\ref{fig:07a1-goods-sale}）、
调货与代销（图~\ref{fig:07a2-goods-transfer}）、送检与处置（图~\ref{fig:07a3-goods-inspect}）。
各迁移及其规约来源见表~\ref{tab:goods-transitions}。

\fig{货品生命周期状态机（八状态总览）}{07a-goods-state}

\fig{货品状态机（一）：售出与预留}{07a1-goods-sale}
\fig{货品状态机（二）：调货与代销}{07a2-goods-transfer}
\fig{货品状态机（三）：送检与处置}{07a3-goods-inspect}

\needspace{8\baselineskip}
\begin{table}[H]
\centering
\small
\caption{货品状态迁移及其规约来源}
\label{tab:goods-transitions}
\begin{tabularx}{\textwidth}{@{}l l X@{}}
\toprule
迁移 & 触发事件 & 规约来源 \\
\midrule
$[*]\rightarrow$在柜 & 入库建档（采购/回收/调货/加工完工/寄售） & UC-02 \\
在柜 $\rightarrow$ 预留 & 客户预留，或预定确定到店 & UC-01 规则③、UC-13 规则⑥ \\
预留 $\rightarrow$ 在柜 & 取消预留（到店看过暂不决定） & UC-13 规则⑥ \\
在柜 $\rightarrow$ 已售 & 全款成交 & UC-01 \\
在柜 $\rightarrow$ 已订 & 收定金、分期留货 & UC-01 规则④ \\
已订 $\rightarrow$ 已售 & 尾款付清并交付出库 & UC-01 规则④ \\
在柜 $\rightarrow$ 已调出 & 调出给合作商户 & UC-20 \\
已调出 $\rightarrow$ 已售 & 同行处售出、鉴赏期届满 & UC-20 规则③ \\
已调出 $\rightarrow$ 在柜 & 退货退回 & UC-20 \\
在柜 $\rightarrow$ 代销中 & 寄卖在同行处 & UC-01 规则⑤ \\
代销中 $\rightarrow$ 已售 & 同行售出、与同行结算 & UC-01 规则⑤ \\
在柜 $\rightarrow$ 已送检 & 送检 & UC-05 \\
预留 $\rightarrow$ 已送检 & 送检 & UC-05 \\
已送检 $\rightarrow$ 在柜 & 证书录入、恢复送检前 & UC-05 \\
已送检 $\rightarrow$ 预留 & 证书录入、恢复送检前 & UC-05 \\
已送检 $\rightarrow$ 待处理 & 检测不合格 & UC-05 扩展 4a \\
已售 $\rightarrow$ 在柜 & 退货（成交后 7 天内、真假/质量） & UC-03 \\
待处理 $\rightarrow$ $[*]$ & 退回供应商或停售 & UC-05 扩展 4a \\
\bottomrule
\end{tabularx}
\end{table}

状态机还体现两条不变式：一是预留、已订期间货品不得再售他人；
二是已调出与代销中的货权不同——已调出货权转移给商户，代销中货权仍归本店。

\subsubsection{出让申请（三状态）}

出让申请的生命周期如图~\ref{fig:07b-apply-state}：供应商提交后进入“待审核”，
老板审核通过则转“已通过”并转入入库流程，驳回则转“已驳回”。
迁移全部来自 UC-14（提交）与 UC-18（审核）。

\fig{出让申请状态机（三状态）}{07b-apply-state}

\subsubsection{加工任务（四状态）}

加工任务的状态机如图~\ref{fig:07c-process-task-state}。它是完工交付申请的组合整体，
“在加工”是提交完工交付申请的前置条件，状态迁移直接由 UC-16、UC-17 驱动。

\fig{加工任务状态机（四状态）}{07c-process-task-state}

\subsubsection{预定记录（四状态）}

预定记录的状态机如图~\ref{fig:07d-booking-state}。额度规则、爽约判定都由状态驱动；
预定本身不锁货，客户确定到店时置为“预留”的是货品状态（见 07a）。

\fig{预定记录状态机（四状态）}{07d-booking-state}

\subsubsection{寄售合同（三状态）}

寄售合同的状态机如图~\ref{fig:07e-consign-state}。登记寄售合同后即“生效”，
提前终止转“已终止”，到期日过仍未终止转“已到期”；
两种结束态下未结算明细对应的货品都退回供应商。

\fig{寄售合同状态机（三状态）}{07e-consign-state}

\subsubsection{分期计划（三状态）}

分期计划的状态机如图~\ref{fig:07f-installment-state}。登记分期计划后为“未还”，
到期日过仍有未还转“逾期”，还清全部转“已结清”；
逾期不是终点，补清欠款后同样进入“已结清”。

\fig{分期计划状态机（三状态）}{07f-installment-state}

% ============================================================
% D 阶段：设计
% ============================================================

% ============================================================
% D 阶段：设计
% ============================================================
\section{D 设计}

设计阶段从分析模型出发“从虚到实”，把分析阶段的类与协作落到具体技术实现。
三层分别选定技术栈（表~\ref{tab:stack}）：
前端 Vue3 + Vite + Element-Plus + Pinia，后端 SpringBoot 3 + MyBatis-Plus，
数据库 MySQL 8.0，外部端微信小程序，柜台电脑端本地部署。
这里按数据层、业务层、表示层三层分别给出设计。

\needspace{8\baselineskip}
\begin{table}[H]
\centering
\small
\caption{三层技术选型}
\label{tab:stack}
\begin{tabularx}{\textwidth}{@{}l X@{}}
\toprule
层 & 技术选型 \\
\midrule
表示层（柜台电脑端） & Vue3 + Vite + Element-Plus + Pinia + Axios，门店本地部署 \\
表示层（外部端） & 微信小程序（客户、供应商、加工师傅、合作商户） \\
业务层 & SpringBoot 3 + MyBatis-Plus \\
数据层 & MySQL 8.0 \\
\bottomrule
\end{tabularx}
\end{table}

\subsection{总体架构（最简架构）}

总体架构如图~\ref{fig:08-architecture}，分三层：表示层（柜台电脑端 + 微信小程序）、
业务层（SpringBoot 3 分层后端）、数据层（MyBatis-Plus 仓储 + MySQL 8.0）。
按最简架构的取舍，能不加的就不加：系统做成单体，不拆微服务；
业务规则只写在业务层一处，货品状态、成交价底线这类跨表规则也由业务层把关，数据库只存取值；
柜台端和小程序调的是同一套后端接口。

\fig{总体架构：表示层 / 业务层 / 数据层}{08-architecture}

\subsection{数据层设计}

数据层把分析实体类映射为 MySQL 8.0 的表：类对表、属性对列、对象对行，
按最简实现满足三范式：每行描述单一对象、无重复组、属性原子（1NF）；
非主属性完全依赖主键、不存在部分依赖（2NF）；非主属性之间不存在传递依赖（3NF）。
原先把客户标签、料子图片、过程图三处多值属性用逗号分隔塞进一个字段，违反第一范式；
现已各自独立成表（\code{t\_customer\_tag}、\code{t\_apply\_image}、\code{t\_delivery\_image}），属性恢复为原子值；
寄售明细（一份合同多件货品）同样拆表（\code{t\_consignment\_item}）。
用户与角色、角色与权限两组多对多各自拆成关联表（\code{t\_user\_role}、\code{t\_role\_permission}）。
供应商、加工师傅、同行商户三类参照数据属性与操作完全一致，合成一张来源方表
\code{t\_source\_party}，以类别列区分，同一类别内名称唯一。
主键一律用代理键（自增 id，唯一且恒定），业务编码单列并加唯一索引；
外键放在“非 1”方（多方），常规业务行用 RESTRICT 禁止物理删除，
多值属性拆表与多对多拆表随主表级联删除（CASCADE），客户介绍人置空（SET NULL）；
外键只保证“指向的记录存在”，来源方类别与单据类型是否匹配、结算单据二选一这类跨表规则由业务层把关；
状态与枚举字段用字符列存储，注释里列全取值。
非规范化只用来微调性能，出现不可调和的性能问题才用。这里只留了一处：
货品表里直接存了证书号，查证书时不用再联表。
全部建表 DDL 见 \code{docs/schema.sql}，共 31 张表（基础参照 6 张 + 业务 19 张 + 拆表 6 张，其中多值属性拆表 4 张、多对多拆表 2 张）。

存储访问在表数据入口、行数据入口、活动记录、数据映射器四种方式中选定\textbf{数据映射器}：
由 MyBatis-Plus 的 Mapper 承担对象与表之间的双向转换，领域层不感知 SQL 与表结构。
这样核心域逻辑与存储逻辑互不污染，日后更换存储也不必改动业务层。

\subsubsection{类-表映射}

分析实体类与数据表的对应关系见表~\ref{tab:class-table}：26 个实体类对应 26 张主表，
另有 6 张拆表（多值属性 4 张：客户标签、出让申请图片、完工过程图、寄售明细；
多对多 2 张：用户-角色、角色-权限），合计 31 张；
基础参照类（来源方、鉴定机构、货品类别、用户、角色、权限）各成一张参照表。
分析模型里的四处组合，落到库上就是四个非空外键：\code{t\_goods.stock\_in\_id}、
\code{t\_certificate.goods\_id}、\code{t\_delivery\_apply.task\_id}、\code{t\_consignment\_item.consignment\_id}；
分析模型中未采用泛化，用户与角色多对多、角色与权限多对多，
角色经 \code{t\_user\_role} 关联，权限经 \code{t\_role\_permission} 关联。

\needspace{8\baselineskip}
\begin{table}[H]
\centering
\small
\caption{分析实体类到数据表的映射（26 个分析类对应 26 张表）}
\label{tab:class-table}
\begin{tabularx}{\textwidth}{@{}l X@{}}
\toprule
分析实体类 & 对应数据表 \\
\midrule
货品（E1） & \code{t\_goods}（货品） \\
货品（E2） & \code{t\_stock\_in}（入库批次） \\
货品（E3） & \code{t\_certificate}（证书） \\
货品（E4） & \code{t\_sale}（成交） \\
货品（E5） & \code{t\_return}（退货） \\
货品（E6） & \code{t\_inspection}（送检） \\
货品（E7） & \code{t\_booking}（预定） \\
货品（E8） & \code{t\_reservation}（预约看货，UC-10） \\
货品（E9） & \code{t\_transfer\_apply}（出让申请，含审核字段） \\
货品（E10） & \code{t\_visit\_appt}（上门预约） \\
货品（E11） & \code{t\_process\_task}（加工任务） \\
货品（E12） & \code{t\_delivery\_apply}（完工交付申请） \\
货品（E13） & \code{t\_payment}（货款） \\
货品（E14） & \code{t\_transfer}（调货） \\
货品（E15） & \code{t\_consignment}（寄售合同） \\
货品（E16） & \code{t\_consignment\_item}（寄售明细，多值拆表） \\
货品（E17） & \code{t\_settlement}（台账结算） \\
货品（E18） & \code{t\_installment}（分期计划） \\
货品（E19） & \code{t\_stock\_log}（库存流水） \\
货品（E20） & \code{t\_customer}（客户） \\
货品（E21） & \code{t\_source\_party}（来源方，类别区分供应商/加工师傅/同行商户） \\
货品（E22） & \code{t\_cert\_org}（鉴定机构） \\
货品（E23） & \code{t\_category}（货品类别） \\
货品（E24） & \code{t\_user}（用户与外部账户） \\
货品（E25） & \code{t\_role}（角色） \\
货品（E26） & \code{t\_permission}（权限） \\
\bottomrule
\end{tabularx}
\end{table}

\subsubsection{分析类名与规约字段列表名的对照}

阶段 A/B 的用例规约「字段列表」用的是业务对象名（如“成交记录”“送检记录”），
分析阶段去掉“记录/数据”后缀，两处所指是同一对象，只是详略程度不同，对照见表~\ref{tab:name-map}。

\needspace{8\baselineskip}
\begin{table}[H]
\centering
\small
\caption{规约字段列表名 / 分析实体类名 / 数据表 的对照}
\label{tab:name-map}
\begin{tabularx}{\textwidth}{@{}p{4.7cm} p{3.2cm} X@{}}
\toprule
规约字段列表名（A/B） & 分析实体类名（C） & 数据表（D） \\
\midrule
成交记录 & 成交 & \code{t\_sale} \\
退货记录 & 退货 & \code{t\_return} \\
送检记录 & 送检 & \code{t\_inspection} \\
预定记录 & 预定 & \code{t\_booking} \\
货款记录 & 货款 & \code{t\_payment} \\
调货记录 & 调货 & \code{t\_transfer} \\
寄售合同 & 寄售合同 & \code{t\_consignment}（明细见 \code{t\_consignment\_item}） \\
台账结算 & 台账结算 & \code{t\_settlement} \\
分期计划 & 分期计划 & \code{t\_installment} \\
库存流水 & 库存流水 & \code{t\_stock\_log} \\
基础数据（来源方/鉴定机构/货品类别） & 同名三个实体类 & 参照表（来源方合一，类别区分） \\
\bottomrule
\end{tabularx}
\end{table}

\needspace{8\baselineskip}
\begin{table}[H]
\centering
\small
\caption{类图上的关系在数据层的落法}
\label{tab:map-rules}
\begin{tabularx}{\textwidth}{@{}p{3.4cm} p{4.0cm} X@{}}
\toprule
分析模型里的 & 落到数据层 & 这里的例子 \\
\midrule
关联 & 外键放在非 1 方（多方） &
\code{t\_sale.goods\_id} 指向 \code{t\_goods}；
\code{t\_certificate.org\_id} 指向 \code{t\_cert\_org} \\
组合 & 外键放在非 1 方，且不为空 &
\code{t\_goods.stock\_in\_id}、\code{t\_certificate.goods\_id}、
\code{t\_delivery\_apply.task\_id}、\code{t\_consignment\_item.consignment\_id} 四列都是 \code{NOT NULL} \\
多值属性 & 拆成独立表，属性恢复原子（1NF） &
\code{t\_customer\_tag}（客户标签）、\code{t\_apply\_image}（出让申请图片）、
\code{t\_delivery\_image}（完工过程图）、\code{t\_consignment\_item}（寄售明细） \\
多对多 & 拆成关联表 &
用户-角色落 \code{t\_user\_role}，角色-权限落 \code{t\_role\_permission}，
两端都是组合外键（非空、级联删除） \\
泛化 & 用超类的主键作为所有类的主键 &
分析模型没有采用泛化；用户与角色多对多、角色与权限多对多，各自拆表；来源方三表合一以类别列区分 \\
主键 & 唯一、恒定，和领域字段分开 &
自增 \code{id} 作主键，业务编码另设一列并加唯一索引，如 \code{t\_goods.goods\_no} \\
\bottomrule
\end{tabularx}
\end{table}

31 张表的整体结构按用途分组列在表~\ref{tab:tables}，
每张表的字段名称、中文含义、数据类型和约束条件在下一节逐张说明。

\subsubsection{表结构清单}

31 张表按用途分组，清单见表~\ref{tab:tables}。

\needspace{6\baselineskip}
\begin{longtable}{@{}p{3.3cm} p{1.9cm} p{5.7cm} p{2.3cm}@{}}
\caption{数据表清单（共 31 张）}\label{tab:tables}\\
\toprule
表名 & 中文名 & 主要字段（列） & 对应用例 \\
\midrule
\endfirsthead
\toprule
表名 & 中文名 & 主要字段（列） & 对应用例 \\
\midrule
\endhead
\bottomrule
\endlastfoot
\code{t\_category} & 货品类别 & name, remark & UC-07 \\
\code{t\_user} & 用户/外部账户 & username, name, password\_hash, status & UC-08/12 \\
\code{t\_role} & 角色 & role\_code, role\_name & UC-12 \\
\code{t\_permission} & 权限 & perm\_code, perm\_name & 全部用例（权限矩阵） \\
\code{t\_source\_party} & 来源方 & party\_type, name, contact\_person, phone & UC-07 \\
\code{t\_cert\_org} & 鉴定机构 & name, phone & UC-07 \\
\code{t\_stock\_in} & 入库批次 & in\_time, receiver\_id & UC-02 \\
\code{t\_goods} & 货品档案 & goods\_no, name, category\_id, material, origin, shape, form, weight\_g, size, cert\_no, image, src\_type, src\_party, src\_doc, cost, min\_price, transfer\_price, in\_time, stock\_in\_id, status & UC-01/02/09/11 \\
\code{t\_customer} & 客户 & customer\_no, name, phone, wechat, source\_channel, introducer\_id, pref\_category, pref\_size, pref\_budget, remark & UC-04 \\
\code{t\_customer\_tag} & 客户标签 & customer\_id, tag & UC-04（多选属性拆表）\\

\code{t\_reservation} & 预约看货 & resv\_no, goods\_id, name, phone, expect\_slot, status & UC-10（状态枚举为设计阶段按接待流程补充）\\
\code{t\_booking} & 预定记录 & booking\_no, goods\_id, name, phone, customer\_id, submit\_time, status & UC-13 \\
\code{t\_inspection} & 送检记录 & insp\_time, goods\_id, org\_id, submitter\_id & UC-05 \\
\code{t\_certificate} & 鉴定证书 & cert\_no, org\_id, insp\_id, goods\_id, conclusion, issue\_date, cert\_qrcode & UC-05 \\
\code{t\_sale} & 成交记录 & sale\_no, goods\_id, sale\_time, price, pay\_method, received, customer\_id, operator\_id, recorded\_at & UC-01/06 \\
\code{t\_return} & 退货记录 & return\_time, sale\_id, reason, refund, operator\_id & UC-03 \\
\code{t\_transfer} & 调货记录 & record\_no, direction, merchant\_id, goods\_id, quantity, status, settled, record\_time & UC-20/23 \\
\code{t\_transfer\_apply} & 出让申请(含审核) & apply\_no, supplier\_id, out\_type, name, quantity, expect\_price, description, submit\_time, status, audit\_opinion, audit\_time, auditor\_id & UC-14/18 \\
\code{t\_apply\_image} & 出让申请图片 & apply\_id, seq, image\_path & UC-14（多图属性拆表）\\

\code{t\_visit\_appt} & 上门预约 & appt\_no, appt\_party, matter, expect\_slot, submit\_time, status & UC-15 \\
\code{t\_process\_task} & 加工任务 & task\_no, goods\_id, craftsman\_id, requirement, labor\_fee, deadline, status & UC-16 \\
\code{t\_delivery\_apply} & 完工交付申请 & apply\_no, task\_id, note, submit\_time, status & UC-17 \\
\code{t\_delivery\_image} & 完工过程图 & apply\_id, stage, image\_path & UC-17（多图属性拆表）\\

\code{t\_payment} & 货款记录 & record\_no, biz\_type, related\_type, related\_no, amount, direction, operator\_id, pay\_time & UC-19 \\
\code{t\_consignment} & 寄售合同 & consign\_no, supplier\_id, sign\_time, expire\_date, settle\_cycle, status & UC-21 \\
\code{t\_consignment\_item} & 寄售明细 & consignment\_id, goods\_id, consign\_price, settle\_status & UC-21/22（多值属性拆表）\\
\code{t\_settlement} & 台账结算 & settle\_no, party\_id, settle\_type, transfer\_id, consign\_item\_id, amount, settle\_method, settle\_time, operator\_id & UC-22/23 \\
\code{t\_installment} & 分期计划 & plan\_no, sale\_id, periods, period\_amount, paid\_amount, next\_due\_date, start\_date, status & UC-24/25 \\
\code{t\_stock\_log} & 库存流水 & goods\_id, from\_status, to\_status, ref\_type, ref\_no, operator\_id, log\_time & UC-26 \\
\code{t\_user\_role} & 用户-角色 & user\_id, role\_id & UC-12（多对多拆表）\\
\code{t\_role\_permission} & 角色-权限 & role\_id, permission\_id & 权限矩阵（多对多拆表）\\
\end{longtable}

\subsubsection{各表字段说明}

下表按案例的写法，逐张列出字段名称、中文描述、数据类型与约束条件。

\needspace{8\baselineskip}
\begin{longtable}{@{}p{3.1cm} p{4.5cm} p{2.4cm} p{4.5cm}@{}}
\caption{货品类别表 \code{t\_category} 的字段说明}\label{tab:f-t-category}\\
\toprule
字段名称 & 中文描述 & 数据类型 & 约束条件 \\
\midrule
\endfirsthead
\toprule
字段名称 & 中文描述 & 数据类型 & 约束条件 \\
\midrule
\endhead
\bottomrule
\endlastfoot
\code{id} & 代理键 & \code{BIGINT UNSIGNED} & 主键、自增、非空 \\
\code{name} & 类别名称（如翡翠/和田玉/碧玉/金丝玉） & \code{VARCHAR(64)} & 唯一、非空 \\
\code{remark} & 备注 & \code{VARCHAR(255)} & 可空 \\
\code{created\_at} & 创建时间 & \code{DATETIME} & 非空 \\
\end{longtable}

\needspace{8\baselineskip}
\begin{longtable}{@{}p{3.1cm} p{4.5cm} p{2.4cm} p{4.5cm}@{}}
\caption{用户与外部账户表 \code{t\_user} 的字段说明}\label{tab:f-t-user}\\
\toprule
字段名称 & 中文描述 & 数据类型 & 约束条件 \\
\midrule
\endfirsthead
\toprule
字段名称 & 中文描述 & 数据类型 & 约束条件 \\
\midrule
\endhead
\bottomrule
\endlastfoot
\code{id} & 代理键 & \code{BIGINT UNSIGNED} & 主键、自增、非空 \\
\code{username} & 用户名（登录名） & \code{VARCHAR(64)} & 唯一、非空 \\
\code{name} & 姓名 & \code{VARCHAR(64)} & 非空 \\
\code{password\_hash} & 口令（加密保存，禁止明文） & \code{VARCHAR(128)} & 非空 \\
\code{status} & 状态：启用|停用 & \code{VARCHAR(8)} & 非空 \\
\code{created\_at} & 创建时间 & \code{DATETIME} & 非空 \\
\code{updated\_at} & 更新时间 & \code{DATETIME} & 非空 \\
\end{longtable}

\needspace{8\baselineskip}
\begin{longtable}{@{}p{3.1cm} p{4.5cm} p{2.4cm} p{4.5cm}@{}}
\caption{角色表 \code{t\_role} 的字段说明}\label{tab:f-t-role}\\
\toprule
字段名称 & 中文描述 & 数据类型 & 约束条件 \\
\midrule
\endfirsthead
\toprule
字段名称 & 中文描述 & 数据类型 & 约束条件 \\
\midrule
\endhead
\bottomrule
\endlastfoot
\code{id} & 代理键 & \code{BIGINT UNSIGNED} & 主键、自增、非空 \\
\code{role\_code} & 角色编码（BOSS/SALE/MERCHANT/SUPPLIER/CRAFTSMAN） & \code{VARCHAR(16)} & 唯一、非空 \\
\code{role\_name} & 角色名：老板|销售|合作商户|供应商|加工师傅 & \code{VARCHAR(32)} & 非空 \\
\code{remark} & 备注 & \code{VARCHAR(255)} & 可空 \\
\code{created\_at} & 创建时间 & \code{DATETIME} & 非空 \\
\end{longtable}

\needspace{8\baselineskip}
\begin{longtable}{@{}p{3.1cm} p{4.5cm} p{2.4cm} p{4.5cm}@{}}
\caption{权限表 \code{t\_permission} 的字段说明}\label{tab:f-t-permission}\\
\toprule
字段名称 & 中文描述 & 数据类型 & 约束条件 \\
\midrule
\endfirsthead
\toprule
字段名称 & 中文描述 & 数据类型 & 约束条件 \\
\midrule
\endhead
\bottomrule
\endlastfoot
\code{id} & 代理键 & \code{BIGINT UNSIGNED} & 主键、自增、非空 \\
\code{perm\_code} & 权限编码（对应用例编号，如 UC-01） & \code{VARCHAR(16)} & 唯一、非空 \\
\code{perm\_name} & 权限名（用例名） & \code{VARCHAR(64)} & 非空 \\
\code{remark} & 备注 & \code{VARCHAR(255)} & 可空 \\
\code{created\_at} & 创建时间 & \code{DATETIME} & 非空 \\
\end{longtable}

\needspace{8\baselineskip}
\begin{longtable}{@{}p{3.1cm} p{4.5cm} p{2.4cm} p{4.5cm}@{}}
\caption{来源方表 \code{t\_source\_party} 的字段说明}\label{tab:f-t-source-party}\\
\toprule
字段名称 & 中文描述 & 数据类型 & 约束条件 \\
\midrule
\endfirsthead
\toprule
字段名称 & 中文描述 & 数据类型 & 约束条件 \\
\midrule
\endhead
\bottomrule
\endlastfoot
\code{id} & 代理键 & \code{BIGINT UNSIGNED} & 主键、自增、非空 \\
\code{party\_type} & 类别：供应商|加工师傅|同行商户 & \code{VARCHAR(16)} & 非空 \\
\code{name} & 名称（公司名/师傅姓名/店名） & \code{VARCHAR(128)} & 类别内唯一、非空 \\
\code{contact\_person} & 联系人（名称为个人时留空） & \code{VARCHAR(64)} & 可空 \\
\code{phone} & 联系电话 & \code{VARCHAR(32)} & 可空 \\
\code{remark} & 备注 & \code{VARCHAR(255)} & 可空 \\
\code{created\_at} & 创建时间 & \code{DATETIME} & 非空 \\
\code{updated\_at} & 更新时间 & \code{DATETIME} & 非空 \\
\end{longtable}

\needspace{8\baselineskip}
\begin{longtable}{@{}p{3.1cm} p{4.5cm} p{2.4cm} p{4.5cm}@{}}
\caption{鉴定机构表 \code{t\_cert\_org} 的字段说明}\label{tab:f-t-cert-org}\\
\toprule
字段名称 & 中文描述 & 数据类型 & 约束条件 \\
\midrule
\endfirsthead
\toprule
字段名称 & 中文描述 & 数据类型 & 约束条件 \\
\midrule
\endhead
\bottomrule
\endlastfoot
\code{id} & 代理键 & \code{BIGINT UNSIGNED} & 主键、自增、非空 \\
\code{name} & 鉴定机构名称 & \code{VARCHAR(128)} & 唯一、非空 \\
\code{phone} & 联系电话 & \code{VARCHAR(32)} & 可空 \\
\code{remark} & 备注 & \code{VARCHAR(255)} & 可空 \\
\code{created\_at} & 创建时间 & \code{DATETIME} & 非空 \\
\code{updated\_at} & 更新时间 & \code{DATETIME} & 非空 \\
\end{longtable}

\needspace{8\baselineskip}
\begin{longtable}{@{}p{3.1cm} p{4.5cm} p{2.4cm} p{4.5cm}@{}}
\caption{入库记录/批次表 \code{t\_stock\_in} 的字段说明}\label{tab:f-t-stock-in}\\
\toprule
字段名称 & 中文描述 & 数据类型 & 约束条件 \\
\midrule
\endfirsthead
\toprule
字段名称 & 中文描述 & 数据类型 & 约束条件 \\
\midrule
\endhead
\bottomrule
\endlastfoot
\code{id} & 代理键 & \code{BIGINT UNSIGNED} & 主键、自增、非空 \\
\code{in\_time} & 入库时间（系统时间） & \code{DATETIME} & 非空 \\
\code{receiver\_id} & 验收人，指向 t\_user & \code{BIGINT UNSIGNED} & 非空、外键（非 1 方）、禁止删除 \\
\code{remark} & 备注 & \code{VARCHAR(255)} & 可空 \\
\code{created\_at} & 创建时间 & \code{DATETIME} & 非空 \\
\end{longtable}

\needspace{8\baselineskip}
\begin{longtable}{@{}p{3.1cm} p{4.5cm} p{2.4cm} p{4.5cm}@{}}
\caption{货品档案表 \code{t\_goods} 的字段说明}\label{tab:f-t-goods}\\
\toprule
字段名称 & 中文描述 & 数据类型 & 约束条件 \\
\midrule
\endfirsthead
\toprule
字段名称 & 中文描述 & 数据类型 & 约束条件 \\
\midrule
\endhead
\bottomrule
\endlastfoot
\code{id} & 代理键 & \code{BIGINT UNSIGNED} & 主键、自增、非空 \\
\code{goods\_no} & 货品编号（条形码，一物一码，系统生成） & \code{VARCHAR(32)} & 唯一、非空 \\
\code{name} & 品名 & \code{VARCHAR(128)} & 非空 \\
\code{category\_id} & 类别，指向 t\_category & \code{BIGINT UNSIGNED} & 非空、外键（非 1 方）、禁止删除 \\
\code{material} & 料种（自定义填写，无法穷举） & \code{VARCHAR(32)} & 非空 \\
\code{origin} & 产地（自定义填写，鉴定证书不判产地，由老板凭眼力判断） & \code{VARCHAR(32)} & 非空 \\
\code{shape} & 器型（自定义填写，无法穷举） & \code{VARCHAR(32)} & 非空 \\
\code{form} & 产状：山料|籽料|山流水料|戈壁料 & \code{VARCHAR(16)} & 非空 \\
\code{weight\_g} & 克重 & \code{DECIMAL(10,2)} & 非空 \\
\code{size} & 尺寸/圈口（如内径56mm） & \code{VARCHAR(64)} & 可空 \\
\code{cert\_no} & 证书号（取最新一张证书的编号，完整证书见 t\_certificate） & \code{VARCHAR(64)} & 可空 \\
\code{image} & 图片（文件路径/URL） & \code{VARCHAR(500)} & 可空 \\
\code{src\_type} & 来源类型：采购|回收|调货|加工完工|寄售 & \code{VARCHAR(16)} & 非空 \\
\code{src\_party} & 来源方（供应商/商户/师傅名） & \code{VARCHAR(128)} & 可空 \\
\code{src\_doc} & 来源单据号 & \code{VARCHAR(64)} & 可空 \\
\code{cost} & 成本（调货调入不记成本，只记往来） & \code{DECIMAL(12,2)} & 可空 \\
\code{min\_price} & 最低售价（成交价不得低于此价） & \code{DECIMAL(12,2)} & 非空 \\
\code{transfer\_price} & 调货价（UC-11 商户端调货价目，在柜时维护） & \code{DECIMAL(12,2)} & 可空 \\
\code{in\_time} & 入库时间（系统时间） & \code{DATETIME} & 非空 \\
\code{stock\_in\_id} & 入库批次，指向 t\_stock\_in（货品清单成员） & \code{BIGINT UNSIGNED} & 非空、外键（非 1 方，组合）、禁止删除 \\
\code{status} & 状态：在柜|已送检|预留|已售|已订|已调出|代销中|待处理 & \code{VARCHAR(16)} & 非空 \\
\code{created\_at} & 创建时间 & \code{DATETIME} & 非空 \\
\code{updated\_at} & 更新时间 & \code{DATETIME} & 非空 \\
\end{longtable}

\needspace{8\baselineskip}
\begin{longtable}{@{}p{3.1cm} p{4.5cm} p{2.4cm} p{4.5cm}@{}}
\caption{客户表 \code{t\_customer} 的字段说明}\label{tab:f-t-customer}\\
\toprule
字段名称 & 中文描述 & 数据类型 & 约束条件 \\
\midrule
\endfirsthead
\toprule
字段名称 & 中文描述 & 数据类型 & 约束条件 \\
\midrule
\endhead
\bottomrule
\endlastfoot
\code{id} & 代理键 & \code{BIGINT UNSIGNED} & 主键、自增、非空 \\
\code{customer\_no} & 客户编号（系统生成） & \code{VARCHAR(32)} & 唯一、非空 \\
\code{name} & 姓名 & \code{VARCHAR(64)} & 非空 \\
\code{phone} & 手机号 & \code{VARCHAR(32)} & 非空、唯一 \\
\code{wechat} & 微信号 & \code{VARCHAR(64)} & 可空 \\
\code{source\_channel} & 来源渠道：闲鱼|微信朋友圈|朋友介绍|线下|直播|老客户复购|其他 & \code{VARCHAR(16)} & 非空 \\
\code{introducer\_id} & 介绍人，指向 t\_customer（来源渠道为朋友介绍时，UC-04） & \code{BIGINT UNSIGNED} & 外键（非 1 方）、置空 \\
\code{pref\_category} & 偏好品类 & \code{VARCHAR(64)} & 可空 \\
\code{pref\_size} & 偏好圈口/尺寸 & \code{VARCHAR(32)} & 可空 \\
\code{pref\_budget} & 偏好预算 & \code{DECIMAL(12,2)} & 可空 \\
\code{remark} & 备注（含跟进记录摘要） & \code{VARCHAR(500)} & 可空 \\
\code{created\_at} & 创建时间 & \code{DATETIME} & 非空 \\
\code{updated\_at} & 更新时间 & \code{DATETIME} & 非空 \\
\end{longtable}

\needspace{8\baselineskip}
\begin{longtable}{@{}p{3.1cm} p{4.5cm} p{2.4cm} p{4.5cm}@{}}
\caption{预约看货表 \code{t\_reservation} 的字段说明}\label{tab:f-t-reservation}\\
\toprule
字段名称 & 中文描述 & 数据类型 & 约束条件 \\
\midrule
\endfirsthead
\toprule
字段名称 & 中文描述 & 数据类型 & 约束条件 \\
\midrule
\endhead
\bottomrule
\endlastfoot
\code{id} & 代理键 & \code{BIGINT UNSIGNED} & 主键、自增、非空 \\
\code{resv\_no} & 预约编号 & \code{VARCHAR(32)} & 唯一、非空 \\
\code{goods\_id} & 货品，指向 t\_goods（预约查看的目标） & \code{BIGINT UNSIGNED} & 非空、外键（非 1 方）、禁止删除 \\
\code{name} & 姓名 & \code{VARCHAR(64)} & 非空 \\
\code{phone} & 手机号 & \code{VARCHAR(32)} & 非空 \\
\code{expect\_slot} & 期望到店时段 & \code{VARCHAR(64)} & 非空 \\
\code{status} & 状态：待确认|已到店|已取消 & \code{VARCHAR(16)} & 非空 \\
\code{created\_at} & 创建时间 & \code{DATETIME} & 非空 \\
\code{updated\_at} & 更新时间 & \code{DATETIME} & 非空 \\
\end{longtable}

\needspace{8\baselineskip}
\begin{longtable}{@{}p{3.1cm} p{4.5cm} p{2.4cm} p{4.5cm}@{}}
\caption{预定记录表 \code{t\_booking} 的字段说明}\label{tab:f-t-booking}\\
\toprule
字段名称 & 中文描述 & 数据类型 & 约束条件 \\
\midrule
\endfirsthead
\toprule
字段名称 & 中文描述 & 数据类型 & 约束条件 \\
\midrule
\endhead
\bottomrule
\endlastfoot
\code{id} & 代理键 & \code{BIGINT UNSIGNED} & 主键、自增、非空 \\
\code{booking\_no} & 预定编号 & \code{VARCHAR(32)} & 唯一、非空 \\
\code{goods\_id} & 货品，指向 t\_goods & \code{BIGINT UNSIGNED} & 非空、外键（非 1 方）、禁止删除 \\
\code{name} & 姓名 & \code{VARCHAR(64)} & 非空 \\
\code{phone} & 手机号（3 次额度按此计） & \code{VARCHAR(32)} & 非空 \\
\code{customer\_id} & 客户，指向 t\_customer（散客可空） & \code{BIGINT UNSIGNED} & 可空、外键（非 1 方）、禁止删除 \\
\code{submit\_time} & 提交时间 & \code{DATETIME} & 非空 \\
\code{status} & 状态：待到店|已达成|已爽约|已取消 & \code{VARCHAR(16)} & 非空 \\
\code{created\_at} & 创建时间 & \code{DATETIME} & 非空 \\
\code{updated\_at} & 更新时间 & \code{DATETIME} & 非空 \\
\end{longtable}

\needspace{8\baselineskip}
\begin{longtable}{@{}p{3.1cm} p{4.5cm} p{2.4cm} p{4.5cm}@{}}
\caption{送检记录表 \code{t\_inspection} 的字段说明}\label{tab:f-t-inspection}\\
\toprule
字段名称 & 中文描述 & 数据类型 & 约束条件 \\
\midrule
\endfirsthead
\toprule
字段名称 & 中文描述 & 数据类型 & 约束条件 \\
\midrule
\endhead
\bottomrule
\endlastfoot
\code{id} & 代理键 & \code{BIGINT UNSIGNED} & 主键、自增、非空 \\
\code{insp\_time} & 送检时间 & \code{DATETIME} & 非空 \\
\code{goods\_id} & 货品，指向 t\_goods & \code{BIGINT UNSIGNED} & 非空、外键（非 1 方）、禁止删除 \\
\code{org\_id} & 送检机构，指向 t\_cert\_org & \code{BIGINT UNSIGNED} & 非空、外键（非 1 方）、禁止删除 \\
\code{submitter\_id} & 送检人，指向 t\_user & \code{BIGINT UNSIGNED} & 非空、外键（非 1 方）、禁止删除 \\
\code{created\_at} & 创建时间 & \code{DATETIME} & 非空 \\
\end{longtable}

\needspace{8\baselineskip}
\begin{longtable}{@{}p{3.1cm} p{4.5cm} p{2.4cm} p{4.5cm}@{}}
\caption{鉴定证书表 \code{t\_certificate} 的字段说明}\label{tab:f-t-certificate}\\
\toprule
字段名称 & 中文描述 & 数据类型 & 约束条件 \\
\midrule
\endfirsthead
\toprule
字段名称 & 中文描述 & 数据类型 & 约束条件 \\
\midrule
\endhead
\bottomrule
\endlastfoot
\code{id} & 代理键 & \code{BIGINT UNSIGNED} & 主键、自增、非空 \\
\code{cert\_no} & 证书编号 & \code{VARCHAR(64)} & 唯一、非空 \\
\code{org\_id} & 鉴定机构，指向 t\_cert\_org & \code{BIGINT UNSIGNED} & 非空、外键（非 1 方）、禁止删除 \\
\code{insp\_id} & 送检记录，指向 t\_inspection（可空） & \code{BIGINT UNSIGNED} & 可空、外键（非 1 方）、禁止删除 \\
\code{goods\_id} & 货品，指向 t\_goods（证书挂到货品） & \code{BIGINT UNSIGNED} & 非空、外键（非 1 方，组合）、禁止删除 \\
\code{conclusion} & 检测结论 & \code{VARCHAR(64)} & 非空 \\
\code{issue\_date} & 签发日期 & \code{DATE} & 非空 \\
\code{cert\_qrcode} & 证书二维码（机构印制，扫码可查） & \code{VARCHAR(500)} & 可空 \\
\code{created\_at} & 创建时间 & \code{DATETIME} & 非空 \\
\end{longtable}

\needspace{8\baselineskip}
\begin{longtable}{@{}p{3.1cm} p{4.5cm} p{2.4cm} p{4.5cm}@{}}
\caption{成交记录表 \code{t\_sale} 的字段说明}\label{tab:f-t-sale}\\
\toprule
字段名称 & 中文描述 & 数据类型 & 约束条件 \\
\midrule
\endfirsthead
\toprule
字段名称 & 中文描述 & 数据类型 & 约束条件 \\
\midrule
\endhead
\bottomrule
\endlastfoot
\code{id} & 代理键 & \code{BIGINT UNSIGNED} & 主键、自增、非空 \\
\code{sale\_no} & 成交编号 & \code{VARCHAR(32)} & 唯一、非空 \\
\code{goods\_id} & 货品，指向 t\_goods & \code{BIGINT UNSIGNED} & 非空、外键（非 1 方）、禁止删除 \\
\code{sale\_time} & 成交发生时间（以系统时间为准） & \code{DATETIME} & 非空 \\
\code{price} & 成交价（不得低于最低售价） & \code{DECIMAL(12,2)} & 非空 \\
\code{pay\_method} & 支付方式：全款|定金 & \code{VARCHAR(8)} & 非空 \\
\code{received} & 已收金额（定金成交时为定金额） & \code{DECIMAL(12,2)} & 非空 \\
\code{customer\_id} & 客户，指向 t\_customer（散客可空） & \code{BIGINT UNSIGNED} & 外键（非 1 方）、禁止删除 \\
\code{operator\_id} & 经手销售，指向 t\_user & \code{BIGINT UNSIGNED} & 非空、外键（非 1 方）、禁止删除 \\
\code{recorded\_at} & 录入时间（与发生时间分别记录） & \code{DATETIME} & 非空 \\
\code{created\_at} & 创建时间 & \code{DATETIME} & 非空 \\
\end{longtable}

\needspace{8\baselineskip}
\begin{longtable}{@{}p{3.1cm} p{4.5cm} p{2.4cm} p{4.5cm}@{}}
\caption{退货记录表 \code{t\_return} 的字段说明}\label{tab:f-t-return}\\
\toprule
字段名称 & 中文描述 & 数据类型 & 约束条件 \\
\midrule
\endfirsthead
\toprule
字段名称 & 中文描述 & 数据类型 & 约束条件 \\
\midrule
\endhead
\bottomrule
\endlastfoot
\code{id} & 代理键 & \code{BIGINT UNSIGNED} & 主键、自增、非空 \\
\code{return\_time} & 退货时间（成交后 7 天内） & \code{DATETIME} & 非空 \\
\code{sale\_id} & 原成交记录，指向 t\_sale & \code{BIGINT UNSIGNED} & 非空、外键（非 1 方）、禁止删除 \\
\code{reason} & 退货原因：真假存疑|质量问题 & \code{VARCHAR(16)} & 非空 \\
\code{refund} & 退款金额 & \code{DECIMAL(12,2)} & 非空 \\
\code{operator\_id} & 经手人，指向 t\_user & \code{BIGINT UNSIGNED} & 非空、外键（非 1 方）、禁止删除 \\
\code{created\_at} & 创建时间 & \code{DATETIME} & 非空 \\
\end{longtable}

\needspace{8\baselineskip}
\begin{longtable}{@{}p{3.1cm} p{4.5cm} p{2.4cm} p{4.5cm}@{}}
\caption{调货记录表 \code{t\_transfer} 的字段说明}\label{tab:f-t-transfer}\\
\toprule
字段名称 & 中文描述 & 数据类型 & 约束条件 \\
\midrule
\endfirsthead
\toprule
字段名称 & 中文描述 & 数据类型 & 约束条件 \\
\midrule
\endhead
\bottomrule
\endlastfoot
\code{id} & 代理键 & \code{BIGINT UNSIGNED} & 主键、自增、非空 \\
\code{record\_no} & 记录编号 & \code{VARCHAR(32)} & 唯一、非空 \\
\code{direction} & 方向：调入|调出 & \code{VARCHAR(8)} & 非空 \\
\code{merchant\_id} & 合作商户，指向 t\_source\_party（类别=同行商户） & \code{BIGINT UNSIGNED} & 非空、外键（非 1 方）、禁止删除 \\
\code{goods\_id} & 货品，指向 t\_goods & \code{BIGINT UNSIGNED} & 非空、外键（非 1 方）、禁止删除 \\
\code{quantity} & 数量 & \code{INT} & 非空 \\
\code{status} & 状态：在调|已售|已退回 & \code{VARCHAR(8)} & 非空 \\
\code{settled} & 已结算：否|是（UC-23 台账结算后置是） & \code{VARCHAR(4)} & 非空 \\
\code{record\_time} & 时间（鉴赏期结算时点见 UC-20 规则③） & \code{DATETIME} & 非空 \\
\code{created\_at} & 创建时间 & \code{DATETIME} & 非空 \\
\code{updated\_at} & 更新时间 & \code{DATETIME} & 非空 \\
\end{longtable}

\needspace{8\baselineskip}
\begin{longtable}{@{}p{3.1cm} p{4.5cm} p{2.4cm} p{4.5cm}@{}}
\caption{出让申请含审核表 \code{t\_transfer\_apply} 的字段说明}\label{tab:f-t-transfer-apply}\\
\toprule
字段名称 & 中文描述 & 数据类型 & 约束条件 \\
\midrule
\endfirsthead
\toprule
字段名称 & 中文描述 & 数据类型 & 约束条件 \\
\midrule
\endhead
\bottomrule
\endlastfoot
\code{id} & 代理键 & \code{BIGINT UNSIGNED} & 主键、自增、非空 \\
\code{apply\_no} & 申请编号 & \code{VARCHAR(32)} & 唯一、非空 \\
\code{supplier\_id} & 供应商，指向 t\_source\_party（类别=供应商，含回流出让方） & \code{BIGINT UNSIGNED} & 非空、外键（非 1 方）、禁止删除 \\
\code{out\_type} & 出让类型：玉料|成品|回流 & \code{VARCHAR(8)} & 非空 \\
\code{name} & 品名 & \code{VARCHAR(128)} & 非空 \\
\code{quantity} & 数量 & \code{INT} & 非空 \\
\code{expect\_price} & 期望出让价 & \code{DECIMAL(12,2)} & 非空 \\
\code{description} & 描述 & \code{VARCHAR(500)} & 可空 \\
\code{submit\_time} & 提交时间 & \code{DATETIME} & 非空 \\
\code{status} & 状态：待审核|已通过|已驳回（状态机见 07b） & \code{VARCHAR(8)} & 非空 \\
\code{audit\_opinion} & 审核意见（UC-18） & \code{VARCHAR(255)} & 可空 \\
\code{audit\_time} & 审核时间（UC-18） & \code{DATETIME} & 可空 \\
\code{auditor\_id} & 审核人，指向 t\_user（UC-18） & \code{BIGINT UNSIGNED} & 外键（非 1 方）、禁止删除 \\
\code{created\_at} & 创建时间 & \code{DATETIME} & 非空 \\
\code{updated\_at} & 更新时间 & \code{DATETIME} & 非空 \\
\end{longtable}

\needspace{8\baselineskip}
\begin{longtable}{@{}p{3.1cm} p{4.5cm} p{2.4cm} p{4.5cm}@{}}
\caption{上门预约表 \code{t\_visit\_appt} 的字段说明}\label{tab:f-t-visit-appt}\\
\toprule
字段名称 & 中文描述 & 数据类型 & 约束条件 \\
\midrule
\endfirsthead
\toprule
字段名称 & 中文描述 & 数据类型 & 约束条件 \\
\midrule
\endhead
\bottomrule
\endlastfoot
\code{id} & 代理键 & \code{BIGINT UNSIGNED} & 主键、自增、非空 \\
\code{appt\_no} & 预约编号 & \code{VARCHAR(32)} & 唯一、非空 \\
\code{appt\_party} & 预约方（供应商/加工师傅） & \code{VARCHAR(128)} & 非空 \\
\code{matter} & 事项（上门交接内容） & \code{VARCHAR(255)} & 非空 \\
\code{expect\_slot} & 期望时段 & \code{VARCHAR(64)} & 非空 \\
\code{submit\_time} & 提交时间 & \code{DATETIME} & 非空 \\
\code{status} & 状态：待上门|已完成|已取消 & \code{VARCHAR(8)} & 非空 \\
\code{transfer\_apply\_id} & 关联出让申请，指向 t\_transfer\_apply（UC-15 规则①，二选一必有其一） & \code{BIGINT UNSIGNED} & 外键（非 1 方）、禁止删除 \\
\code{delivery\_apply\_id} & 关联完工交付申请，指向 t\_delivery\_apply（UC-15 规则①，二选一必有其一） & \code{BIGINT UNSIGNED} & 外键（非 1 方）、禁止删除 \\
\code{created\_at} & 创建时间 & \code{DATETIME} & 非空 \\
\code{updated\_at} & 更新时间 & \code{DATETIME} & 非空 \\
\end{longtable}

\needspace{8\baselineskip}
\begin{longtable}{@{}p{3.1cm} p{4.5cm} p{2.4cm} p{4.5cm}@{}}
\caption{加工任务表 \code{t\_process\_task} 的字段说明}\label{tab:f-t-process-task}\\
\toprule
字段名称 & 中文描述 & 数据类型 & 约束条件 \\
\midrule
\endfirsthead
\toprule
字段名称 & 中文描述 & 数据类型 & 约束条件 \\
\midrule
\endhead
\bottomrule
\endlastfoot
\code{id} & 代理键 & \code{BIGINT UNSIGNED} & 主键、自增、非空 \\
\code{task\_no} & 任务编号 & \code{VARCHAR(32)} & 唯一、非空 \\
\code{goods\_id} & 货品，指向 t\_goods（加工对象） & \code{BIGINT UNSIGNED} & 非空、外键（非 1 方）、禁止删除 \\
\code{craftsman\_id} & 加工师傅，指向 t\_source\_party（类别=加工师傅，分派对象） & \code{BIGINT UNSIGNED} & 非空、外键（非 1 方）、禁止删除 \\
\code{requirement} & 加工要求 & \code{VARCHAR(500)} & 非空 \\
\code{labor\_fee} & 工费 & \code{DECIMAL(12,2)} & 非空 \\
\code{deadline} & 期限 & \code{DATE} & 非空 \\
\code{status} & 状态：待加工|在加工|待交付|已交付 & \code{VARCHAR(8)} & 非空 \\
\code{created\_at} & 创建时间 & \code{DATETIME} & 非空 \\
\code{updated\_at} & 更新时间 & \code{DATETIME} & 非空 \\
\end{longtable}

\needspace{8\baselineskip}
\begin{longtable}{@{}p{3.1cm} p{4.5cm} p{2.4cm} p{4.5cm}@{}}
\caption{完工交付申请表 \code{t\_delivery\_apply} 的字段说明}\label{tab:f-t-delivery-apply}\\
\toprule
字段名称 & 中文描述 & 数据类型 & 约束条件 \\
\midrule
\endfirsthead
\toprule
字段名称 & 中文描述 & 数据类型 & 约束条件 \\
\midrule
\endhead
\bottomrule
\endlastfoot
\code{id} & 代理键 & \code{BIGINT UNSIGNED} & 主键、自增、非空 \\
\code{apply\_no} & 申请编号 & \code{VARCHAR(32)} & 唯一、非空 \\
\code{task\_id} & 任务编号，指向 t\_process\_task & \code{BIGINT UNSIGNED} & 非空、外键（非 1 方，组合）、禁止删除 \\
\code{note} & 完工说明 & \code{VARCHAR(500)} & 可空 \\
\code{submit\_time} & 提交时间 & \code{DATETIME} & 非空 \\
\code{status} & 状态：待确认|已确认（确认后安排入库） & \code{VARCHAR(8)} & 非空 \\
\code{created\_at} & 创建时间 & \code{DATETIME} & 非空 \\
\code{updated\_at} & 更新时间 & \code{DATETIME} & 非空 \\
\end{longtable}

\needspace{8\baselineskip}
\begin{longtable}{@{}p{3.1cm} p{4.5cm} p{2.4cm} p{4.5cm}@{}}
\caption{货款记录表 \code{t\_payment} 的字段说明}\label{tab:f-t-payment}\\
\toprule
字段名称 & 中文描述 & 数据类型 & 约束条件 \\
\midrule
\endfirsthead
\toprule
字段名称 & 中文描述 & 数据类型 & 约束条件 \\
\midrule
\endhead
\bottomrule
\endlastfoot
\code{id} & 代理键 & \code{BIGINT UNSIGNED} & 主键、自增、非空 \\
\code{record\_no} & 记录编号 & \code{VARCHAR(32)} & 唯一、非空 \\
\code{biz\_type} & 业务类型：采购货款|调货货款|加工工费|寄售货款 & \code{VARCHAR(16)} & 非空 \\
\code{related\_type} & 关联单据类型（入库批次/调货记录/加工任务） & \code{VARCHAR(32)} & 非空 \\
\code{related\_no} & 关联单据编号（逻辑引用，不建硬外键） & \code{VARCHAR(32)} & 非空 \\
\code{amount} & 金额 & \code{DECIMAL(12,2)} & 非空 \\
\code{direction} & 收付方向：支付|收取 & \code{VARCHAR(8)} & 非空 \\
\code{operator\_id} & 经手人，指向 t\_user & \code{BIGINT UNSIGNED} & 非空、外键（非 1 方）、禁止删除 \\
\code{pay\_time} & 时间 & \code{DATETIME} & 非空 \\
\code{created\_at} & 创建时间 & \code{DATETIME} & 非空 \\
\end{longtable}

\needspace{8\baselineskip}
\begin{longtable}{@{}p{3.1cm} p{4.5cm} p{2.4cm} p{4.5cm}@{}}
\caption{客户标签表 \code{t\_customer\_tag} 的字段说明}\label{tab:f-t-customer-tag}\\
\toprule
字段名称 & 中文描述 & 数据类型 & 约束条件 \\
\midrule
\endfirsthead
\toprule
字段名称 & 中文描述 & 数据类型 & 约束条件 \\
\midrule
\endhead
\bottomrule
\endlastfoot
\code{id} & 代理键 & \code{BIGINT UNSIGNED} & 主键、自增、非空 \\
\code{customer\_id} & 客户，指向 t\_customer & \code{BIGINT UNSIGNED} & 非空、外键（非 1 方）、随主表级联删除 \\
\code{tag} & 标签：大客户|懂行|爽快|要还价 & \code{VARCHAR(16)} & 非空 \\
\code{created\_at} & 创建时间 & \code{DATETIME} & 非空 \\
\end{longtable}

\needspace{8\baselineskip}
\begin{longtable}{@{}p{3.1cm} p{4.5cm} p{2.4cm} p{4.5cm}@{}}
\caption{出让申请图片表 \code{t\_apply\_image} 的字段说明}\label{tab:f-t-apply-image}\\
\toprule
字段名称 & 中文描述 & 数据类型 & 约束条件 \\
\midrule
\endfirsthead
\toprule
字段名称 & 中文描述 & 数据类型 & 约束条件 \\
\midrule
\endhead
\bottomrule
\endlastfoot
\code{id} & 代理键 & \code{BIGINT UNSIGNED} & 主键、自增、非空 \\
\code{apply\_id} & 出让申请，指向 t\_transfer\_apply & \code{BIGINT UNSIGNED} & 非空、外键（非 1 方）、随主表级联删除 \\
\code{seq} & 显示顺序 & \code{INT} & 非空 \\
\code{image\_path} & 图片路径 & \code{VARCHAR(255)} & 非空 \\
\end{longtable}

\needspace{8\baselineskip}
\begin{longtable}{@{}p{3.1cm} p{4.5cm} p{2.4cm} p{4.5cm}@{}}
\caption{完工交付过程图表 \code{t\_delivery\_image} 的字段说明}\label{tab:f-t-delivery-image}\\
\toprule
字段名称 & 中文描述 & 数据类型 & 约束条件 \\
\midrule
\endfirsthead
\toprule
字段名称 & 中文描述 & 数据类型 & 约束条件 \\
\midrule
\endhead
\bottomrule
\endlastfoot
\code{id} & 代理键 & \code{BIGINT UNSIGNED} & 主键、自增、非空 \\
\code{apply\_id} & 完工交付申请，指向 t\_delivery\_apply & \code{BIGINT UNSIGNED} & 非空、外键（非 1 方）、随主表级联删除 \\
\code{stage} & 阶段：加工前|加工中|完工 & \code{VARCHAR(16)} & 非空 \\
\code{image\_path} & 图片路径 & \code{VARCHAR(255)} & 非空 \\
\end{longtable}

\needspace{8\baselineskip}
\begin{longtable}{@{}p{3.1cm} p{4.5cm} p{2.4cm} p{4.5cm}@{}}
\caption{寄售合同表 \code{t\_consignment} 的字段说明}\label{tab:f-t-consignment}\\
\toprule
字段名称 & 中文描述 & 数据类型 & 约束条件 \\
\midrule
\endfirsthead
\toprule
字段名称 & 中文描述 & 数据类型 & 约束条件 \\
\midrule
\endhead
\bottomrule
\endlastfoot
\code{id} & 代理键 & \code{BIGINT UNSIGNED} & 主键、自增、非空 \\
\code{consign\_no} & 合同编号 & \code{VARCHAR(32)} & 唯一、非空 \\
\code{supplier\_id} & 寄售供应商，指向 t\_source\_party（类别=供应商） & \code{BIGINT UNSIGNED} & 非空、外键（非 1 方）、禁止删除 \\
\code{sign\_time} & 签订时间 & \code{DATETIME} & 非空 \\
\code{expire\_date} & 到期日（到期未终止即置已到期） & \code{DATE} & 非空 \\
\code{settle\_cycle} & 结算方式：售出即结|按月结算 & \code{VARCHAR(16)} & 非空 \\
\code{status} & 状态：生效|已终止|已到期（状态机见 07e） & \code{VARCHAR(8)} & 非空 \\
\code{remark} & 备注 & \code{VARCHAR(255)} & 可空 \\
\code{created\_at} & 创建时间 & \code{DATETIME} & 非空 \\
\code{updated\_at} & 更新时间 & \code{DATETIME} & 非空 \\
\end{longtable}

\needspace{8\baselineskip}
\begin{longtable}{@{}p{3.1cm} p{4.5cm} p{2.4cm} p{4.5cm}@{}}
\caption{寄售明细表 \code{t\_consignment\_item} 的字段说明}\label{tab:f-t-consignment-item}\\
\toprule
字段名称 & 中文描述 & 数据类型 & 约束条件 \\
\midrule
\endfirsthead
\toprule
字段名称 & 中文描述 & 数据类型 & 约束条件 \\
\midrule
\endhead
\bottomrule
\endlastfoot
\code{id} & 代理键 & \code{BIGINT UNSIGNED} & 主键、自增、非空 \\
\code{consignment\_id} & 寄售合同，指向 t\_consignment & \code{BIGINT UNSIGNED} & 非空、组合外键（非 1 方）、随主表级联删除 \\
\code{goods\_id} & 货品，指向 t\_goods & \code{BIGINT UNSIGNED} & 非空、索引、外键（非 1 方）、禁止删除 \\
\code{consign\_price} & 寄售价（售出后按此价与供应商结算） & \code{DECIMAL(12,2)} & 非空 \\
\code{settle\_status} & 结算状态：未结算|已结算（UC-22） & \code{VARCHAR(8)} & 非空 \\
\code{created\_at} & 创建时间 & \code{DATETIME} & 非空 \\
\end{longtable}

\needspace{8\baselineskip}
\begin{longtable}{@{}p{3.1cm} p{4.5cm} p{2.4cm} p{4.5cm}@{}}
\caption{台账结算表 \code{t\_settlement} 的字段说明}\label{tab:f-t-settlement}\\
\toprule
字段名称 & 中文描述 & 数据类型 & 约束条件 \\
\midrule
\endfirsthead
\toprule
字段名称 & 中文描述 & 数据类型 & 约束条件 \\
\midrule
\endhead
\bottomrule
\endlastfoot
\code{id} & 代理键 & \code{BIGINT UNSIGNED} & 主键、自增、非空 \\
\code{settle\_no} & 结算编号 & \code{VARCHAR(32)} & 唯一、非空 \\
\code{party\_id} & 结算往来方，指向 t\_source\_party & \code{BIGINT UNSIGNED} & 非空、外键（非 1 方）、禁止删除 \\
\code{settle\_type} & 结算类型：调货|寄售 & \code{VARCHAR(8)} & 非空 \\
\code{transfer\_id} & 调货记录，指向 t\_transfer（结算类型=调货时非空） & \code{BIGINT UNSIGNED} & 外键（非 1 方）、禁止删除 \\
\code{consign\_item\_id} & 寄售明细，指向 t\_consignment\_item（结算类型=寄售时非空） & \code{BIGINT UNSIGNED} & 外键（非 1 方）、禁止删除 \\
\code{amount} & 结算金额 & \code{DECIMAL(12,2)} & 非空 \\
\code{settle\_method} & 结算方式：现金|转账|冲抵 & \code{VARCHAR(16)} & 非空 \\
\code{settle\_time} & 结算时间 & \code{DATETIME} & 非空 \\
\code{operator\_id} & 经手人，指向 t\_user & \code{BIGINT UNSIGNED} & 非空、外键（非 1 方）、禁止删除 \\
\code{created\_at} & 创建时间 & \code{DATETIME} & 非空 \\
\end{longtable}

\needspace{8\baselineskip}
\begin{longtable}{@{}p{3.1cm} p{4.5cm} p{2.4cm} p{4.5cm}@{}}
\caption{分期计划表 \code{t\_installment} 的字段说明}\label{tab:f-t-installment}\\
\toprule
字段名称 & 中文描述 & 数据类型 & 约束条件 \\
\midrule
\endfirsthead
\toprule
字段名称 & 中文描述 & 数据类型 & 约束条件 \\
\midrule
\endhead
\bottomrule
\endlastfoot
\code{id} & 代理键 & \code{BIGINT UNSIGNED} & 主键、自增、非空 \\
\code{plan\_no} & 计划编号 & \code{VARCHAR(32)} & 唯一、非空 \\
\code{sale\_id} & 成交记录，指向 t\_sale & \code{BIGINT UNSIGNED} & 非空、外键（非 1 方）、禁止删除 \\
\code{periods} & 分期期数 & \code{INT} & 非空 \\
\code{period\_amount} & 每期应还金额 & \code{DECIMAL(12,2)} & 非空 \\
\code{paid\_amount} & 已还金额 & \code{DECIMAL(12,2)} & 非空 \\
\code{next\_due\_date} & 下一期到期日（超过未结清即逾期） & \code{DATE} & 非空 \\
\code{start\_date} & 计划生效日 & \code{DATE} & 非空 \\
\code{status} & 状态：未还|逾期|已结清（状态机见 07f） & \code{VARCHAR(8)} & 非空 \\
\code{created\_at} & 创建时间 & \code{DATETIME} & 非空 \\
\code{updated\_at} & 更新时间 & \code{DATETIME} & 非空 \\
\end{longtable}

\needspace{8\baselineskip}
\begin{longtable}{@{}p{3.1cm} p{4.5cm} p{2.4cm} p{4.5cm}@{}}
\caption{库存流水表 \code{t\_stock\_log} 的字段说明}\label{tab:f-t-stock-log}\\
\toprule
字段名称 & 中文描述 & 数据类型 & 约束条件 \\
\midrule
\endfirsthead
\toprule
字段名称 & 中文描述 & 数据类型 & 约束条件 \\
\midrule
\endhead
\bottomrule
\endlastfoot
\code{id} & 代理键 & \code{BIGINT UNSIGNED} & 主键、自增、非空 \\
\code{goods\_id} & 货品，指向 t\_goods & \code{BIGINT UNSIGNED} & 非空、外键（非 1 方）、禁止删除 \\
\code{from\_status} & 变更前状态（入库为第一条流水时为空） & \code{VARCHAR(16)} & 可空 \\
\code{to\_status} & 变更后状态（取值同 t\_goods.status） & \code{VARCHAR(16)} & 非空 \\
\code{ref\_type} & 关联单据类型（成交/退货/调货记录/加工任务/送检记录/入库批次/寄售合同） & \code{VARCHAR(16)} & 可空 \\
\code{ref\_no} & 关联单据编号（逻辑引用，同 t\_payment 口径） & \code{VARCHAR(32)} & 可空 \\
\code{operator\_id} & 经手人，指向 t\_user & \code{BIGINT UNSIGNED} & 非空、外键（非 1 方）、禁止删除 \\
\code{log\_time} & 流水时间（系统时间） & \code{DATETIME} & 非空 \\
\code{created\_at} & 创建时间 & \code{DATETIME} & 非空 \\
\end{longtable}

\needspace{8\baselineskip}
\begin{longtable}{@{}p{3.1cm} p{4.5cm} p{2.4cm} p{4.5cm}@{}}
\caption{用户-角色表 \code{t\_user\_role} 的字段说明}\label{tab:f-t-user-role}\\
\toprule
字段名称 & 中文描述 & 数据类型 & 约束条件 \\
\midrule
\endfirsthead
\toprule
字段名称 & 中文描述 & 数据类型 & 约束条件 \\
\midrule
\endhead
\bottomrule
\endlastfoot
\code{id} & 代理键 & \code{BIGINT UNSIGNED} & 主键、自增、非空 \\
\code{user\_id} & 用户，指向 t\_user & \code{BIGINT UNSIGNED} & 非空、组合外键（非 1 方）、随主表级联删除 \\
\code{role\_id} & 角色，指向 t\_role & \code{BIGINT UNSIGNED} & 非空、组合外键（非 1 方）、随主表级联删除 \\
\end{longtable}

\needspace{8\baselineskip}
\begin{longtable}{@{}p{3.1cm} p{4.5cm} p{2.4cm} p{4.5cm}@{}}
\caption{角色-权限表 \code{t\_role\_permission} 的字段说明}\label{tab:f-t-role-permission}\\
\toprule
字段名称 & 中文描述 & 数据类型 & 约束条件 \\
\midrule
\endfirsthead
\toprule
字段名称 & 中文描述 & 数据类型 & 约束条件 \\
\midrule
\endhead
\bottomrule
\endlastfoot
\code{id} & 代理键 & \code{BIGINT UNSIGNED} & 主键、自增、非空 \\
\code{role\_id} & 角色，指向 t\_role & \code{BIGINT UNSIGNED} & 非空、组合外键（非 1 方）、随主表级联删除 \\
\code{permission\_id} & 权限，指向 t\_permission & \code{BIGINT UNSIGNED} & 非空、组合外键（非 1 方）、随主表级联删除 \\
\end{longtable}

\subsubsection{建表 SQL}

31 张表的完整建表语句（含字段注释、唯一索引、外键约束）见 \code{docs/schema.sql}。
设计要点：建库统一 \code{utf8mb4} 字符集；每张表以自增 \code{id} 为代理键；
业务编码列（goods\_no、sale\_no 等）加唯一索引；外键集中在非 1 方，
常规业务行 \code{ON DELETE RESTRICT}，多值属性拆表与多对多拆表 \code{ON DELETE CASCADE}，客户介绍人 \code{ON DELETE SET NULL}；
金额用 \code{DECIMAL(12,2)}。
来源方三表合一，类别列区分、类别内名称唯一（唯一键 类别+名称）；客户手机号唯一（同一手机号不重复建档）；
台账结算一单一据，调货记录与寄售明细两条可空外键二选一，由业务层把关；
货款记录关联的单据跨入库、调货、加工等多类，用 \code{biz\_type + related\_type + related\_no} 定位；
库存流水只增不改不删，货品每次状态变更必写一条。

\subsection{业务层设计}

业务层用 SpringBoot 3 + MyBatis-Plus 按 Controller / Service / Mapper 分层：
Controller 接收请求与参数校验，Service 承载全部业务规则（规则只写在这一层），
Mapper 通过 MyBatis-Plus 封装为仓储（数据映射器）访问数据层，屏蔽 SQL 细节。

类图给出属性、关系与操作的定义；序列图给出操作实现中调用其他类的操作；
状态机图给出操作实现中的迁移与触发事件；警戒条件已在 B 段规约业务规则中写明，进入/退出动作由 Service 执行。
服务的操作从序列图和状态机推出来，不是从属性推出来。
服务划分见表~\ref{tab:services}：服务是控制类落到业务层的载体，一个服务承接同域的若干控制类
（如 CollabService 覆盖出让、加工、预约、预定 7 个控制类），职责与 C 阶段控制类一致。

\needspace{8\baselineskip}
\begin{table}[H]
\centering
\small
\caption{业务层服务划分与职责}
\label{tab:services}
\begin{tabularx}{\textwidth}{@{}l X l@{}}
\toprule
服务 & 职责 & 主要用例 \\
\midrule
GoodsService & 入库建档、货品状态变更、按金额加权周转、库存流水写入与查询、基础数据维护 & UC-01/02/05/07/20/06/26 \\
SaleService & 成交登记、退货与统计冲减、分期计划与还款 & UC-01/03/24/25 \\
CustomerService & 客户档案维护与检索、偏好 & UC-04 \\
TransferService & 调货往来与调货库存维护 & UC-20 \\
PaymentService & 货款与工费收付记录 & UC-19 \\
ConsignService & 寄售合同登记与寄售货款结算 & UC-21/22 \\
LedgerService & 调货台账结算 & UC-23 \\
CollabService & 货品浏览、调货价目、出让申请/审核、加工任务/完工、预约、预定 & UC-09/10/11/13/14/15/16/17/18 \\
StatsService & 月度销售额、销售客户数、渠道分布 & UC-06 \\
AuthService & 登录鉴权、角色与外部账户开通、权限矩阵校验 & UC-08/12 \\
\bottomrule
\end{tabularx}
\end{table}

以登记成交为例，请求自表现层经 Axios 到达 \code{SaleController}，
\code{SaleService} 校验货品处于“在柜”、成交价不低于最低售价，
再分别经 \code{GoodsMapper} 更新货品状态、经 \code{SaleMapper} 落成交记录，
最后返回登记结果。

跨表不变式由业务层统一把关，全部写在 Service 一层：
①货品每次状态变更必须写一条库存流水，流水写入后不可修改、不可删除；
②一张台账结算单只对应一张业务单据（调货记录或寄售明细），结算后单据置已结算，已结算单据不再出现在结算清单；
③一份成交至多一份分期计划，计划总额等于成交价，还款不超过未还余额；
④一件货品同一时间只挂一份未结清的寄售明细；
⑤用户至少拥有一个角色、角色至少拥有一条权限，权限与用例一一对应；
⑥来源方类别与引用它的单据类型匹配（出让申请引用类别=供应商的来源方，加工任务引用类别=加工师傅，调货记录引用类别=同行商户）。

\subsection{表示层设计}

表示层分两端：柜台电脑端（门店本地，服务老板与销售）与微信小程序（外部协同端，
服务客户、供应商、加工师傅、合作商户，操作步骤尽量少）。
功能模块划分见表~\ref{tab:modules}。

\needspace{8\baselineskip}
\begin{table}[H]
\centering
\small
\caption{表示层功能模块划分}
\label{tab:modules}
\begin{tabularx}{\textwidth}{@{}l X@{}}
\toprule
端 / 模块 & 功能 \\
\midrule
柜台电脑端 · 老板管理端 & 货品入库、统计报表、基础数据维护、外部账户开通、出让申请审核、货款登记、调货记录、寄售合同（UC-21）、寄售与台账结算（UC-22/23）、库存流水（UC-26） \\
柜台电脑端 · 销售柜台端 & 登记成交、登记退货、客户维护、送检与证书登记、分期计划与还款（UC-24/25） \\
微信小程序 · 客户 & 浏览货品（UC-09）、预约看货（UC-10）、预定商品（UC-13） \\
微信小程序 · 供应商 & 提交出让申请（UC-14）、预约上门时间（UC-15） \\
微信小程序 · 加工师傅 & 预约上门时间（UC-15）、查看加工任务（UC-16）、提交完工交付申请（UC-17） \\
微信小程序 · 合作商户 & 查看调货价目（UC-11） \\
\bottomrule
\end{tabularx}
\end{table}

柜台电脑端用 Vue3 组件化组织，Axios 统一封装对后端接口的调用；
微信小程序各角色登录后只暴露与该角色相关的少数字段与按钮，降低操作步骤。
两端共用同一套后端接口，角色与权限由业务层 AuthService 统一控制。
表示层的组件按角色细分，分析阶段的边界类按外部系统划分（柜台电脑端、微信小程序共 2 个），
两者粒度不同：一个边界类对应某一端里的多个角色组件。
状态机也要落到界面上：货品处在不同状态时可点的按钮不同，
例如已售的货品不再显示“登记成交”，在检的货品显示“录入证书”。

\subsubsection{页面状态机与组件可用}

页面自己的状态决定组件什么时候可用。以两个代表页面为例给出页面状态机：
成交登记页（录入型）与调货结算页（操作型），如图~\ref{fig:11a-page-state-deal}
与图~\ref{fig:11b-page-state-transfer}；其余页面按同一模板给出。

\fig{设计状态机图：成交登记页的页面状态}{11a-page-state-deal}

\fig{设计状态机图：调货结算页的页面状态}{11b-page-state-transfer}

页面状态的迁移全部由用户事件与校验结果驱动，与 C 阶段状态机一致：
成交登记页的“校验通过”即 UC-01 的全部校验规则，调货结算页的“创建结算单成功”即 UC-23 的结算规则。
每个页面状态下组件的可用与禁用见表~\ref{tab:ui-deal} 与表~\ref{tab:ui-transfer}，
禁用的组件置灰显示而不是隐藏，操作者能看出功能存在但不满足条件。

\needspace{8\baselineskip}
\begin{table}[H]
\centering
\small
\caption{成交登记页组件可用表}
\label{tab:ui-deal}
\begin{tabularx}{\textwidth}{@{}l l X@{}}
\toprule
组件 & 类型 & 可用条件（其余状态置灰） \\
\midrule
货品选择列表 & 输入 & 待选货品、信息填写状态下可用；检索范围仅“在柜”货品 \\
成交价输入 & 输入 & 信息填写状态下可用；不高于最低售价时提交按钮置灰并提示 \\
客户选择 & 输入 & 信息填写状态下可用；可空（散客） \\
支付方式选择 & 输入 & 信息填写状态下可用；全款或定金 \\
分期入口按钮 & 按钮 & 信息填写状态下、成交金额达到分期门槛时可用，进入 UC-24 \\
提交按钮 & 按钮 & 信息填写状态下且必填项完整时可用；校验中、已登记状态下置灰 \\
返回按钮 & 按钮 & 任意状态可用 \\
\bottomrule
\end{tabularx}
\end{table}

\needspace{8\baselineskip}
\begin{table}[H]
\centering
\small
\caption{调货结算页组件可用表}
\label{tab:ui-transfer}
\begin{tabularx}{\textwidth}{@{}l l X@{}}
\toprule
组件 & 类型 & 可用条件（其余状态置灰） \\
\midrule
商户选择 & 输入 & 待选商户状态下可用；检索范围为基础数据中类别为同行商户的来源方 \\
记录清单 & 显示 & 记录清单状态下载入该商户未结算的调货记录，只读 \\
记录勾选 & 输入 & 记录清单、清单确认状态下可用；仅“已结算=否”的记录可选 \\
结算金额 & 显示 & 只读，系统按勾选记录汇总，不可录入 \\
结算方式选择 & 输入 & 清单确认状态下可用；现金、转账或冲抵 \\
结算按钮 & 按钮 & 清单确认状态下且至少勾选一条记录时可用；结算中、已结算状态下置灰 \\
返回按钮 & 按钮 & 任意状态可用 \\
\bottomrule
\end{tabularx}
\end{table}

\subsubsection{权限矩阵与角色控制}

用户与角色多对多、角色与权限多对多（见 D 段数据层）：
一个账号可以同时具有多个角色，一个角色聚合多条权限，权限与用例一一对应。
登录时 AuthService 按 \code{t\_user\_role} 取出角色集合，再经 \code{t\_role\_permission}
汇成权限集合，菜单按权限生成、按钮按权限矩阵置灰。
五个角色对全部用例的授权见表~\ref{tab:perm}（√ 为可执行，— 为无权限）。

\needspace{8\baselineskip}
\begin{table}[H]
\centering
\small
\caption{角色功能权限矩阵}
\label{tab:perm}
\begin{tabularx}{\textwidth}{@{}X c c c c c@{}}
\toprule
用例（权限） & 老板 & 销售 & 合作商户 & 供应商 & 加工师傅 \\
\midrule
UC-01 登记成交 / UC-03 登记退货 & √ & √ & — & — & — \\
UC-02 登记货品入库 & √ & — & — & — & — \\
UC-04 维护客户信息 & √ & √ & — & — & — \\
UC-05 登记送检与证书 & √ & √ & — & — & — \\
UC-06 查看销售与周转统计 & √ & — & — & — & — \\
UC-07 维护基础数据 & √ & — & — & — & — \\
UC-09 浏览货品 / UC-10 预约看货 / UC-13 预定商品 & — & — & — & — & — \\
UC-11 查看调货价格 & — & — & √ & — & — \\
UC-12 开通外部账户 & √ & — & — & — & — \\
UC-14 提交货品出让申请 / UC-15 预约上门 & — & — & — & √ & √ \\
UC-16 查看加工任务 / UC-17 提交完工交付申请 & — & — & — & — & √ \\
UC-18 审核出让申请 & √ & — & — & — & — \\
UC-19 登记货款记录 & √ & — & — & — & — \\
UC-20 维护调货记录 & √ & 只读 & — & — & — \\
UC-21 登记寄售合同 / UC-22 结算寄售货款 & √ & — & — & — & — \\
UC-23 结算调货台账 & √ & — & — & — & — \\
UC-24 登记分期计划 / UC-25 记录分期还款 & √ & √ & — & — & — \\
UC-26 查看库存流水 & √ & 只读 & — & — & — \\
\bottomrule
\end{tabularx}
\end{table}

注：老板在销售端用例行的 √ 为代执行权限，不改变 B 段规约的主执行者。

矩阵落在库上就是 \code{t\_role\_permission} 的种子数据：每行一个角色对一条权限的授权，
老板持有全部权限，外部三角色只持各自提交与查询类权限。
权限校验在业务层 AuthService 收口，表示层按同一矩阵控制菜单与按钮的置灰，
两层读的是同一份数据，不会出现界面放行而接口拒绝之外的口径差。

\end{document}
